\href{https://lilianweng.github.io/}{Lil\textquotesingle Log}

\begin{itemize}
\tightlist
\item
  \textbar{}
\end{itemize}

\begin{itemize}
\tightlist
\item
  \href{https://lilianweng.github.io/}{{Posts}}
\item
  \href{https://lilianweng.github.io/archives}{{Archive}}
\item
  \href{https://lilianweng.github.io/search/}{{Search}}
\item
  \href{https://lilianweng.github.io/tags/}{{Tags}}
\item
  \href{https://lilianweng.github.io/faq}{{FAQ}}
\end{itemize}

\section{Policy Gradient Algorithms}\label{policy-gradient-algorithms}

Date: April 8, 2018 \textbar{} Estimated Reading Time: 52 min \textbar{} Author: Lilian Weng

{Table of Contents}

\begin{itemize}
\tightlist
\item
  \hyperref[what-is-policy-gradient]{What is Policy Gradient}

  \begin{itemize}
  \tightlist
  \item
    \hyperref[notations]{Notations}
  \item
    \hyperref[policy-gradient]{Policy Gradient}
  \item
    \hyperref[policy-gradient-theorem]{Policy Gradient Theorem}
  \item
    \hyperref[proof-of-policy-gradient-theorem]{Proof of Policy Gradient Theorem}
  \end{itemize}
\item
  \hyperref[policy-gradient-algorithms]{Policy Gradient Algorithms}

  \begin{itemize}
  \tightlist
  \item
    \hyperref[reinforce]{REINFORCE}
  \item
    \hyperref[actor-critic]{Actor-Critic}
  \item
    \hyperref[off-policy-policy-gradient]{Off-Policy Policy Gradient}
  \item
    \hyperref[a3c]{A3C}
  \item
    \hyperref[a2c]{A2C}
  \item
    \hyperref[dpg]{DPG}
  \item
    \hyperref[ddpg]{DDPG}
  \item
    \hyperref[d4pg]{D4PG}
  \item
    \hyperref[maddpg]{MADDPG}
  \item
    \hyperref[trpo]{TRPO}
  \item
    \hyperref[ppo]{PPO}
  \item
    \hyperref[ppg]{PPG}
  \item
    \hyperref[acer]{ACER}
  \item
    \hyperref[actkr]{ACTKR}
  \item
    \hyperref[sac]{SAC}
  \item
    \hyperref[sac-with-automatically-adjusted-temperature]{SAC with Automatically Adjusted Temperature}
  \item
    \hyperref[td3]{TD3}
  \item
    \hyperref[svpg]{SVPG}
  \item
    \hyperref[impala]{IMPALA}
  \end{itemize}
\item
  \hyperref[quick-summary]{Quick Summary}
\item
  \hyperref[references]{References}
\end{itemize}

{{[}Updated on 2018-06-30: add two new policy gradient methods, \hyperref[sac]{SAC} and \hyperref[d4pg]{D4PG}.{]}}\\
{{[}Updated on 2018-09-30: add a new policy gradient method, \hyperref[td3]{TD3}.{]}}\\
{{[}Updated on 2019-02-09: add \hyperref[sac-with-automatically-adjusted-temperature]{SAC with automatically adjusted temperature}{]}.}\\
{{[}Updated on 2019-06-26: Thanks to Chanseok, we have a version of this post in \href{https://talkingaboutme.tistory.com/entry/RL-Policy-Gradient-Algorithms}{Korean}{]}.}\\
{{[}Updated on 2019-09-12: add a new policy gradient method \hyperref[svpg]{SVPG}.{]}}\\
{{[}Updated on 2019-12-22: add a new policy gradient method \hyperref[impala]{IMPALA}.{]}}\\
{{[}Updated on 2020-10-15: add a new policy gradient method \hyperref[ppg]{PPG} \& some new discussion in \hyperref[ppo]{PPO}.{]}}\\
{{[}Updated on 2021-09-19: Thanks to Wenhao \& 爱吃猫的鱼, we have this post in \href{https://tomaxent.com/2019/04/14/\%E7\%AD\%96\%E7\%95\%A5\%E6\%A2\%AF\%E5\%BA\%A6\%E6\%96\%B9\%E6\%B3\%95/}{Chinese1} \& \href{https://paperexplained.cn/articles/article/detail/31/}{Chinese2}{]}.}

\section{\texorpdfstring{What is Policy Gradient\hyperref[what-is-policy-gradient]{\#}}{What is Policy Gradient\#}}\label{what-is-policy-gradient}

Policy gradient is an approach to solve reinforcement learning problems. If you haven't looked into the field of reinforcement learning, please first read the section \href{https://lilianweng.github.io/posts/2018-02-19-rl-overview/\#key-concepts}{``A (Long) Peek into Reinforcement Learning » Key Concepts''} for the problem definition and key concepts.

\subsection{\texorpdfstring{Notations\hyperref[notations]{\#}}{Notations\#}}\label{notations}

Here is a list of notations to help you read through equations in the post easily.

\begin{longtable}[]{@{}ll@{}}
\toprule\noalign{}
Symbol & Meaning \\
\midrule\noalign{}
\endhead
\bottomrule\noalign{}
\endlastfoot
\$s \textbackslash in \textbackslash mathcal\{S\}\$ & States. \\
\$a \textbackslash in \textbackslash mathcal\{A\}\$ & Actions. \\
\$r \textbackslash in \textbackslash mathcal\{R\}\$ & Rewards. \\
\$S\_t, A\_t, R\_t\$ & State, action, and reward at time step \$t\$ of one trajectory. I may occasionally use \$s\_t, a\_t, r\_t\$ as well. \\
\$\textbackslash gamma\$ & Discount factor; penalty to uncertainty of future rewards; \$0\textless\textbackslash gamma \textbackslash leq 1\$. \\
\$G\_t\$ & Return; or discounted future reward; \$G\_t = \textbackslash sum\_\{k=0\}\^{}\{\textbackslash infty\} \textbackslash gamma\^{}k R\_\{t+k+1\}\$. \\
\$P(s', r \textbackslash vert s, a)\$ & Transition probability of getting to the next state \$s'\$ from the current state \$s\$ with action \$a\$ and reward \$r\$. \\
\$\textbackslash pi(a \textbackslash vert s)\$ & Stochastic policy (agent behavior strategy); \$\textbackslash pi\_\textbackslash theta(.)\$ is a policy parameterized by \$\textbackslash theta\$. \\
\$\textbackslash mu(s)\$ & Deterministic policy; we can also label this as \$\textbackslash pi(s)\$, but using a different letter gives better distinction so that we can easily tell when the policy is stochastic or deterministic without further explanation. Either \$\textbackslash pi\$ or \$\textbackslash mu\$ is what a reinforcement learning algorithm aims to learn. \\
\$V(s)\$ & State-value function measures the expected return of state \$s\$; \$V\_w(.)\$ is a value function parameterized by \$w\$. \\
\$V\^{}\textbackslash pi(s)\$ & The value of state \$s\$ when we follow a policy \$\textbackslash pi\$; \$V\^{}\textbackslash pi (s) = \textbackslash mathbb\{E\}\_\{a\textbackslash sim \textbackslash pi\} {[}G\_t \textbackslash vert S\_t = s{]}\$. \\
\$Q(s, a)\$ & Action-value function is similar to \$V(s)\$, but it assesses the expected return of a pair of state and action \$(s, a)\$; \$Q\_w(.)\$ is a action value function parameterized by \$w\$. \\
\$Q\^{}\textbackslash pi(s, a)\$ & Similar to \$V\^{}\textbackslash pi(.)\$, the value of (state, action) pair when we follow a policy \$\textbackslash pi\$; \$Q\^{}\textbackslash pi(s, a) = \textbackslash mathbb\{E\}\_\{a\textbackslash sim \textbackslash pi\} {[}G\_t \textbackslash vert S\_t = s, A\_t = a{]}\$. \\
\$A(s, a)\$ & Advantage function, \$A(s, a) = Q(s, a) - V(s)\$; it can be considered as another version of Q-value with lower variance by taking the state-value off as the baseline. \\
\end{longtable}

\subsection{\texorpdfstring{Policy Gradient\hyperref[policy-gradient]{\#}}{Policy Gradient\#}}\label{policy-gradient}

The goal of reinforcement learning is to find an optimal behavior strategy for the agent to obtain optimal rewards. The \textbf{policy gradient} methods target at modeling and optimizing the policy directly. The policy is usually modeled with a parameterized function respect to \$\textbackslash theta\$, \$\textbackslash pi\_\textbackslash theta(a \textbackslash vert s)\$. The value of the reward (objective) function depends on this policy and then various algorithms can be applied to optimize \$\textbackslash theta\$ for the best reward.

The reward function is defined as:

\$\$ J(\textbackslash theta) = \textbackslash sum\_\{s \textbackslash in \textbackslash mathcal\{S\}\} d\^{}\textbackslash pi(s) V\^{}\textbackslash pi(s) = \textbackslash sum\_\{s \textbackslash in \textbackslash mathcal\{S\}\} d\^{}\textbackslash pi(s) \textbackslash sum\_\{a \textbackslash in \textbackslash mathcal\{A\}\} \textbackslash pi\_\textbackslash theta(a \textbackslash vert s) Q\^{}\textbackslash pi(s, a) \$\$

where \$d\^{}\textbackslash pi(s)\$ is the stationary distribution of Markov chain for \$\textbackslash pi\_\textbackslash theta\$ (on-policy state distribution under \$\textbackslash pi\$). For simplicity, the parameter \$\textbackslash theta\$ would be omitted for the policy \$\textbackslash pi\_\textbackslash theta\$ when the policy is present in the subscript of other functions; for example, \$d\^{}\{\textbackslash pi\}\$ and \$Q\^{}\textbackslash pi\$ should be \$d\^{}\{\textbackslash pi\_\textbackslash theta\}\$ and \$Q\^{}\{\textbackslash pi\_\textbackslash theta\}\$ if written in full.

Imagine that you can travel along the Markov chain's states forever, and eventually, as the time progresses, the probability of you ending up with one state becomes unchanged --- this is the stationary probability for \$\textbackslash pi\_\textbackslash theta\$. \$d\^{}\textbackslash pi(s) = \textbackslash lim\_\{t \textbackslash to \textbackslash infty\} P(s\_t = s \textbackslash vert s\_0, \textbackslash pi\_\textbackslash theta)\$ is the probability that \$s\_t=s\$ when starting from \$s\_0\$ and following policy \$\textbackslash pi\_\textbackslash theta\$ for t steps. Actually, the existence of the stationary distribution of Markov chain is one main reason for why PageRank algorithm works. If you want to read more, check \href{https://jeremykun.com/2015/04/06/markov-chain-monte-carlo-without-all-the-bullshit/}{this}.

It is natural to expect policy-based methods are more useful in the continuous space. Because there is an infinite number of actions and (or) states to estimate the values for and hence value-based approaches are way too expensive computationally in the continuous space. For example, in \href{https://lilianweng.github.io/posts/2018-02-19-rl-overview/\#policy-iteration}{generalized policy iteration}, the policy improvement step \$\textbackslash arg\textbackslash max\_\{a \textbackslash in \textbackslash mathcal\{A\}\} Q\^{}\textbackslash pi(s, a)\$ requires a full scan of the action space, suffering from the \href{https://en.wikipedia.org/wiki/Curse_of_dimensionality}{curse of dimensionality}.

Using \emph{gradient ascent}, we can move \$\textbackslash theta\$ toward the direction suggested by the gradient \$\textbackslash nabla\_\textbackslash theta J(\textbackslash theta)\$ to find the best \$\textbackslash theta\$ for \$\textbackslash pi\_\textbackslash theta\$ that produces the highest return.

\subsection{\texorpdfstring{Policy Gradient Theorem\hyperref[policy-gradient-theorem]{\#}}{Policy Gradient Theorem\#}}\label{policy-gradient-theorem}

Computing the gradient \$\textbackslash nabla\_\textbackslash theta J(\textbackslash theta)\$ is tricky because it depends on both the action selection (directly determined by \$\textbackslash pi\_\textbackslash theta\$) and the stationary distribution of states following the target selection behavior (indirectly determined by \$\textbackslash pi\_\textbackslash theta\$). Given that the environment is generally unknown, it is difficult to estimate the effect on the state distribution by a policy update.

Luckily, the \textbf{policy gradient theorem} comes to save the world! Woohoo! It provides a nice reformation of the derivative of the objective function to not involve the derivative of the state distribution \$d\^{}\textbackslash pi(.)\$ and simplify the gradient computation \$\textbackslash nabla\_\textbackslash theta J(\textbackslash theta)\$ a lot.

\$\$ \textbackslash begin\{aligned\} \textbackslash nabla\_\textbackslash theta J(\textbackslash theta) \&= \textbackslash nabla\_\textbackslash theta \textbackslash sum\_\{s \textbackslash in \textbackslash mathcal\{S\}\} d\^{}\textbackslash pi(s) \textbackslash sum\_\{a \textbackslash in \textbackslash mathcal\{A\}\} Q\^{}\textbackslash pi(s, a) \textbackslash pi\_\textbackslash theta(a \textbackslash vert s) \textbackslash\textbackslash{} \&\textbackslash propto \textbackslash sum\_\{s \textbackslash in \textbackslash mathcal\{S\}\} d\^{}\textbackslash pi(s) \textbackslash sum\_\{a \textbackslash in \textbackslash mathcal\{A\}\} Q\^{}\textbackslash pi(s, a) \textbackslash nabla\_\textbackslash theta \textbackslash pi\_\textbackslash theta(a \textbackslash vert s) \textbackslash end\{aligned\} \$\$

\subsection{\texorpdfstring{Proof of Policy Gradient Theorem\hyperref[proof-of-policy-gradient-theorem]{\#}}{Proof of Policy Gradient Theorem\#}}\label{proof-of-policy-gradient-theorem}

This session is pretty dense, as it is the time for us to go through the proof (\href{http://incompleteideas.net/book/bookdraft2017nov5.pdf}{Sutton \& Barto, 2017}; Sec. 13.1) and figure out why the policy gradient theorem is correct.

We first start with the derivative of the state value function:

\$\$ \textbackslash begin\{aligned\} \& \textbackslash nabla\_\textbackslash theta V\^{}\textbackslash pi(s) \textbackslash\textbackslash{} =\& \textbackslash nabla\_\textbackslash theta \textbackslash Big(\textbackslash sum\_\{a \textbackslash in \textbackslash mathcal\{A\}\} \textbackslash pi\_\textbackslash theta(a \textbackslash vert s)Q\^{}\textbackslash pi(s, a) \textbackslash Big) \& \textbackslash\textbackslash{} =\& \textbackslash sum\_\{a \textbackslash in \textbackslash mathcal\{A\}\} \textbackslash Big( \textbackslash nabla\_\textbackslash theta \textbackslash pi\_\textbackslash theta(a \textbackslash vert s)Q\^{}\textbackslash pi(s, a) + \textbackslash pi\_\textbackslash theta(a \textbackslash vert s) \textbackslash color\{red\}\{\textbackslash nabla\_\textbackslash theta Q\^{}\textbackslash pi(s, a)\} \textbackslash Big) \& \textbackslash scriptstyle\{\textbackslash text\{; Derivative product rule.\}\} \textbackslash\textbackslash{} =\& \textbackslash sum\_\{a \textbackslash in \textbackslash mathcal\{A\}\} \textbackslash Big( \textbackslash nabla\_\textbackslash theta \textbackslash pi\_\textbackslash theta(a \textbackslash vert s)Q\^{}\textbackslash pi(s, a) + \textbackslash pi\_\textbackslash theta(a \textbackslash vert s) \textbackslash color\{red\}\{\textbackslash nabla\_\textbackslash theta \textbackslash sum\_\{s\textquotesingle, r\} P(s\textquotesingle,r \textbackslash vert s,a)(r + V\^{}\textbackslash pi(s\textquotesingle))\} \textbackslash Big) \& \textbackslash scriptstyle\{\textbackslash text\{; Extend \} Q\^{}\textbackslash pi \textbackslash text\{ with future state value.\}\} \textbackslash\textbackslash{} =\& \textbackslash sum\_\{a \textbackslash in \textbackslash mathcal\{A\}\} \textbackslash Big( \textbackslash nabla\_\textbackslash theta \textbackslash pi\_\textbackslash theta(a \textbackslash vert s)Q\^{}\textbackslash pi(s, a) + \textbackslash pi\_\textbackslash theta(a \textbackslash vert s) \textbackslash color\{red\}\{\textbackslash sum\_\{s\textquotesingle, r\} P(s\textquotesingle,r \textbackslash vert s,a) \textbackslash nabla\_\textbackslash theta V\^{}\textbackslash pi(s\textquotesingle)\} \textbackslash Big) \& \textbackslash scriptstyle\{P(s\textquotesingle,r \textbackslash vert s,a) \textbackslash text\{ or \} r \textbackslash text\{ is not a func of \}\textbackslash theta\}\textbackslash\textbackslash{} =\& \textbackslash sum\_\{a \textbackslash in \textbackslash mathcal\{A\}\} \textbackslash Big( \textbackslash nabla\_\textbackslash theta \textbackslash pi\_\textbackslash theta(a \textbackslash vert s)Q\^{}\textbackslash pi(s, a) + \textbackslash pi\_\textbackslash theta(a \textbackslash vert s) \textbackslash color\{red\}\{\textbackslash sum\_\{s\textquotesingle\} P(s\textquotesingle{} \textbackslash vert s,a) \textbackslash nabla\_\textbackslash theta V\^{}\textbackslash pi(s\textquotesingle)\} \textbackslash Big) \& \textbackslash scriptstyle\{\textbackslash text\{; Because \} P(s\textquotesingle{} \textbackslash vert s, a) = \textbackslash sum\_r P(s\textquotesingle, r \textbackslash vert s, a)\} \textbackslash end\{aligned\} \$\$

Now we have:

\$\$ \textbackslash color\{red\}\{\textbackslash nabla\_\textbackslash theta V\^{}\textbackslash pi(s)\} = \textbackslash sum\_\{a \textbackslash in \textbackslash mathcal\{A\}\} \textbackslash Big( \textbackslash nabla\_\textbackslash theta \textbackslash pi\_\textbackslash theta(a \textbackslash vert s)Q\^{}\textbackslash pi(s, a) + \textbackslash pi\_\textbackslash theta(a \textbackslash vert s) \textbackslash sum\_\{s\textquotesingle\} P(s\textquotesingle{} \textbackslash vert s,a) \textbackslash color\{red\}\{\textbackslash nabla\_\textbackslash theta V\^{}\textbackslash pi(s\textquotesingle)\} \textbackslash Big) \$\$

This equation has a nice recursive form (see the red parts!) and the future state value function \$V\^{}\textbackslash pi(s')\$ can be repeated unrolled by following the same equation.

Let's consider the following visitation sequence and label the probability of transitioning from state s to state x with policy \$\textbackslash pi\_\textbackslash theta\$ after k step as \$\textbackslash rho\^{}\textbackslash pi(s \textbackslash to x, k)\$.

\$\$ s \textbackslash xrightarrow{[}{]}\{a \textbackslash sim \textbackslash pi\_\textbackslash theta(.\textbackslash vert s)\} s\textquotesingle{} \textbackslash xrightarrow{[}{]}\{a \textbackslash sim \textbackslash pi\_\textbackslash theta(.\textbackslash vert s\textquotesingle)\} s\textquotesingle\textquotesingle{} \textbackslash xrightarrow{[}{]}\{a \textbackslash sim \textbackslash pi\_\textbackslash theta(.\textbackslash vert s\textquotesingle\textquotesingle)\} \textbackslash dots \$\$

\begin{itemize}
\tightlist
\item
  When k = 0: \$\textbackslash rho\^{}\textbackslash pi(s \textbackslash to s, k=0) = 1\$.
\item
  When k = 1, we scan through all possible actions and sum up the transition probabilities to the target state: \$\textbackslash rho\^{}\textbackslash pi(s \textbackslash to s', k=1) = \textbackslash sum\_a \textbackslash pi\_\textbackslash theta(a \textbackslash vert s) P(s' \textbackslash vert s, a)\$.
\item
  Imagine that the goal is to go from state s to x after k+1 steps while following policy \$\textbackslash pi\_\textbackslash theta\$. We can first travel from s to a middle point s' (any state can be a middle point, \$s' \textbackslash in \textbackslash mathcal\{S\}\$) after k steps and then go to the final state x during the last step. In this way, we are able to update the visitation probability recursively: \$\textbackslash rho\^{}\textbackslash pi(s \textbackslash to x, k+1) = \textbackslash sum\_\{s'\} \textbackslash rho\^{}\textbackslash pi(s \textbackslash to s', k) \textbackslash rho\^{}\textbackslash pi(s' \textbackslash to x, 1)\$.
\end{itemize}

Then we go back to unroll the recursive representation of \$\textbackslash nabla\_\textbackslash theta V\^{}\textbackslash pi(s)\$! Let \$\textbackslash phi(s) = \textbackslash sum\_\{a \textbackslash in \textbackslash mathcal\{A\}\} \textbackslash nabla\_\textbackslash theta \textbackslash pi\_\textbackslash theta(a \textbackslash vert s)Q\^{}\textbackslash pi(s, a)\$ to simplify the maths. If we keep on extending \$\textbackslash nabla\_\textbackslash theta V\^{}\textbackslash pi(.)\$ infinitely, it is easy to find out that we can transition from the starting state s to any state after any number of steps in this unrolling process and by summing up all the visitation probabilities, we get \$\textbackslash nabla\_\textbackslash theta V\^{}\textbackslash pi(s)\$!

\$\$ \textbackslash begin\{aligned\} \& \textbackslash color\{red\}\{\textbackslash nabla\_\textbackslash theta V\^{}\textbackslash pi(s)\} \textbackslash\textbackslash{} =\& \textbackslash phi(s) + \textbackslash sum\_a \textbackslash pi\_\textbackslash theta(a \textbackslash vert s) \textbackslash sum\_\{s\textquotesingle\} P(s\textquotesingle{} \textbackslash vert s,a) \textbackslash color\{red\}\{\textbackslash nabla\_\textbackslash theta V\^{}\textbackslash pi(s\textquotesingle)\} \textbackslash\textbackslash{} =\& \textbackslash phi(s) + \textbackslash sum\_\{s\textquotesingle\} \textbackslash sum\_a \textbackslash pi\_\textbackslash theta(a \textbackslash vert s) P(s\textquotesingle{} \textbackslash vert s,a) \textbackslash color\{red\}\{\textbackslash nabla\_\textbackslash theta V\^{}\textbackslash pi(s\textquotesingle)\} \textbackslash\textbackslash{} =\& \textbackslash phi(s) + \textbackslash sum\_\{s\textquotesingle\} \textbackslash rho\^{}\textbackslash pi(s \textbackslash to s\textquotesingle, 1) \textbackslash color\{red\}\{\textbackslash nabla\_\textbackslash theta V\^{}\textbackslash pi(s\textquotesingle)\} \textbackslash\textbackslash{} =\& \textbackslash phi(s) + \textbackslash sum\_\{s\textquotesingle\} \textbackslash rho\^{}\textbackslash pi(s \textbackslash to s\textquotesingle, 1) \textbackslash color\{red\}\{\textbackslash sum\_\{a \textbackslash in \textbackslash mathcal\{A\}\} \textbackslash Big( \textbackslash nabla\_\textbackslash theta \textbackslash pi\_\textbackslash theta(a \textbackslash vert s\textquotesingle)Q\^{}\textbackslash pi(s\textquotesingle, a) + \textbackslash pi\_\textbackslash theta(a \textbackslash vert s\textquotesingle) \textbackslash sum\_\{s\textquotesingle\} P(s\textquotesingle\textquotesingle{} \textbackslash vert s\textquotesingle,a) \textbackslash nabla\_\textbackslash theta V\^{}\textbackslash pi(s\textquotesingle\textquotesingle) \textbackslash Big)\} \textbackslash\textbackslash{} =\& \textbackslash phi(s) + \textbackslash sum\_\{s\textquotesingle\} \textbackslash rho\^{}\textbackslash pi(s \textbackslash to s\textquotesingle, 1) \textbackslash color\{red\}\{{[} \textbackslash phi(s\textquotesingle) + \textbackslash sum\_\{s\textquotesingle\textquotesingle\} \textbackslash rho\^{}\textbackslash pi(s\textquotesingle{} \textbackslash to s\textquotesingle\textquotesingle, 1) \textbackslash nabla\_\textbackslash theta V\^{}\textbackslash pi(s\textquotesingle\textquotesingle){]}\} \textbackslash\textbackslash{} =\& \textbackslash phi(s) + \textbackslash sum\_\{s\textquotesingle\} \textbackslash rho\^{}\textbackslash pi(s \textbackslash to s\textquotesingle, 1) \textbackslash phi(s\textquotesingle) + \textbackslash sum\_\{s\textquotesingle\textquotesingle\} \textbackslash rho\^{}\textbackslash pi(s \textbackslash to s\textquotesingle\textquotesingle, 2)\textbackslash color\{red\}\{\textbackslash nabla\_\textbackslash theta V\^{}\textbackslash pi(s\textquotesingle\textquotesingle)\} \textbackslash scriptstyle\{\textbackslash text\{ ; Consider \}s\textquotesingle\textbackslash text\{ as the middle point for \}s \textbackslash to s\textquotesingle\textquotesingle\}\textbackslash\textbackslash{} =\& \textbackslash phi(s) + \textbackslash sum\_\{s\textquotesingle\} \textbackslash rho\^{}\textbackslash pi(s \textbackslash to s\textquotesingle, 1) \textbackslash phi(s\textquotesingle) + \textbackslash sum\_\{s\textquotesingle\textquotesingle\} \textbackslash rho\^{}\textbackslash pi(s \textbackslash to s\textquotesingle\textquotesingle, 2)\textbackslash phi(s\textquotesingle\textquotesingle) + \textbackslash sum\_\{s\textquotesingle\textquotesingle\textquotesingle\} \textbackslash rho\^{}\textbackslash pi(s \textbackslash to s\textquotesingle\textquotesingle\textquotesingle, 3)\textbackslash color\{red\}\{\textbackslash nabla\_\textbackslash theta V\^{}\textbackslash pi(s\textquotesingle\textquotesingle\textquotesingle)\} \textbackslash\textbackslash{} =\& \textbackslash dots \textbackslash scriptstyle\{\textbackslash text\{; Repeatedly unrolling the part of \}\textbackslash nabla\_\textbackslash theta V\^{}\textbackslash pi(.)\} \textbackslash\textbackslash{} =\& \textbackslash sum\_\{x\textbackslash in\textbackslash mathcal\{S\}\}\textbackslash sum\_\{k=0\}\^{}\textbackslash infty \textbackslash rho\^{}\textbackslash pi(s \textbackslash to x, k) \textbackslash phi(x) \textbackslash end\{aligned\} \$\$

The nice rewriting above allows us to exclude the derivative of Q-value function, \$\textbackslash nabla\_\textbackslash theta Q\^{}\textbackslash pi(s, a)\$. By plugging it into the objective function \$J(\textbackslash theta)\$, we are getting the following:

\$\$ \textbackslash begin\{aligned\} \textbackslash nabla\_\textbackslash theta J(\textbackslash theta) \&= \textbackslash nabla\_\textbackslash theta V\^{}\textbackslash pi(s\_0) \& \textbackslash scriptstyle\{\textbackslash text\{; Starting from a random state \} s\_0\} \textbackslash\textbackslash{} \&= \textbackslash sum\_\{s\}\textbackslash color\{blue\}\{\textbackslash sum\_\{k=0\}\^{}\textbackslash infty \textbackslash rho\^{}\textbackslash pi(s\_0 \textbackslash to s, k)\} \textbackslash phi(s) \&\textbackslash scriptstyle\{\textbackslash text\{; Let \}\textbackslash color\{blue\}\{\textbackslash eta(s) = \textbackslash sum\_\{k=0\}\^{}\textbackslash infty \textbackslash rho\^{}\textbackslash pi(s\_0 \textbackslash to s, k)\}\} \textbackslash\textbackslash{} \&= \textbackslash sum\_\{s\}\textbackslash eta(s) \textbackslash phi(s) \& \textbackslash\textbackslash{} \&= \textbackslash Big( \{\textbackslash sum\_s \textbackslash eta(s)\} \textbackslash Big)\textbackslash sum\_\{s\}\textbackslash frac\{\textbackslash eta(s)\}\{\textbackslash sum\_s \textbackslash eta(s)\} \textbackslash phi(s) \& \textbackslash scriptstyle\{\textbackslash text\{; Normalize \} \textbackslash eta(s), s\textbackslash in\textbackslash mathcal\{S\} \textbackslash text\{ to be a probability distribution.\}\}\textbackslash\textbackslash{} \&\textbackslash propto \textbackslash sum\_s \textbackslash frac\{\textbackslash eta(s)\}\{\textbackslash sum\_s \textbackslash eta(s)\} \textbackslash phi(s) \& \textbackslash scriptstyle\{\textbackslash sum\_s \textbackslash eta(s)\textbackslash text\{ is a constant\}\} \textbackslash\textbackslash{} \&= \textbackslash sum\_s d\^{}\textbackslash pi(s) \textbackslash sum\_a \textbackslash nabla\_\textbackslash theta \textbackslash pi\_\textbackslash theta(a \textbackslash vert s)Q\^{}\textbackslash pi(s, a) \& \textbackslash scriptstyle\{d\^{}\textbackslash pi(s) = \textbackslash frac\{\textbackslash eta(s)\}\{\textbackslash sum\_s \textbackslash eta(s)\}\textbackslash text\{ is stationary distribution.\}\} \textbackslash end\{aligned\} \$\$

In the episodic case, the constant of proportionality (\$\textbackslash sum\_s \textbackslash eta(s)\$) is the average length of an episode; in the continuing case, it is 1 (\href{http://incompleteideas.net/book/bookdraft2017nov5.pdf}{Sutton \& Barto, 2017}; Sec. 13.2). The gradient can be further written as:

\$\$ \textbackslash begin\{aligned\} \textbackslash nabla\_\textbackslash theta J(\textbackslash theta) \&\textbackslash propto \textbackslash sum\_\{s \textbackslash in \textbackslash mathcal\{S\}\} d\^{}\textbackslash pi(s) \textbackslash sum\_\{a \textbackslash in \textbackslash mathcal\{A\}\} Q\^{}\textbackslash pi(s, a) \textbackslash nabla\_\textbackslash theta \textbackslash pi\_\textbackslash theta(a \textbackslash vert s) \&\textbackslash\textbackslash{} \&= \textbackslash sum\_\{s \textbackslash in \textbackslash mathcal\{S\}\} d\^{}\textbackslash pi(s) \textbackslash sum\_\{a \textbackslash in \textbackslash mathcal\{A\}\} \textbackslash pi\_\textbackslash theta(a \textbackslash vert s) Q\^{}\textbackslash pi(s, a) \textbackslash frac\{\textbackslash nabla\_\textbackslash theta \textbackslash pi\_\textbackslash theta(a \textbackslash vert s)\}\{\textbackslash pi\_\textbackslash theta(a \textbackslash vert s)\} \&\textbackslash\textbackslash{} \&= \textbackslash mathbb\{E\}\_\textbackslash pi {[}Q\^{}\textbackslash pi(s, a) \textbackslash nabla\_\textbackslash theta \textbackslash ln \textbackslash pi\_\textbackslash theta(a \textbackslash vert s){]} \& \textbackslash scriptstyle\{\textbackslash text\{; Because \} (\textbackslash ln x)\textquotesingle{} = 1/x\} \textbackslash end\{aligned\} \$\$

Where \$\textbackslash mathbb\{E\}\_\textbackslash pi\$ refers to \$\textbackslash mathbb\{E\}\_\{s \textbackslash sim d\_\textbackslash pi, a \textbackslash sim \textbackslash pi\_\textbackslash theta\}\$ when both state and action distributions follow the policy \$\textbackslash pi\_\textbackslash theta\$ (on policy).

The policy gradient theorem lays the theoretical foundation for various policy gradient algorithms. This vanilla policy gradient update has no bias but high variance. Many following algorithms were proposed to reduce the variance while keeping the bias unchanged.

\$\$ \textbackslash nabla\_\textbackslash theta J(\textbackslash theta) = \textbackslash mathbb\{E\}\_\textbackslash pi {[}Q\^{}\textbackslash pi(s, a) \textbackslash nabla\_\textbackslash theta \textbackslash ln \textbackslash pi\_\textbackslash theta(a \textbackslash vert s){]} \$\$

Here is a nice summary of a general form of policy gradient methods borrowed from the \href{https://arxiv.org/pdf/1506.02438.pdf}{GAE} (general advantage estimation) paper (\href{https://arxiv.org/abs/1506.02438}{Schulman et al., 2016}) and this \href{https://danieltakeshi.github.io/2017/04/02/notes-on-the-generalized-advantage-estimation-paper/}{post} thoroughly discussed several components in GAE , highly recommended.

\begin{figure}
\centering
\includegraphics{general_form_policy_gradient.png}
\caption{A general form of policy gradient methods. (Image source: \href{https://arxiv.org/abs/1506.02438}{Schulman et al., 2016})}
\end{figure}

\section{\texorpdfstring{Policy Gradient Algorithms\hyperref[policy-gradient-algorithms]{\#}}{Policy Gradient Algorithms\#}}\label{policy-gradient-algorithms}

Tons of policy gradient algorithms have been proposed during recent years and there is no way for me to exhaust them. I'm introducing some of them that I happened to know and read about.

\subsection{\texorpdfstring{REINFORCE\hyperref[reinforce]{\#}}{REINFORCE\#}}\label{reinforce}

\textbf{REINFORCE} (Monte-Carlo policy gradient) relies on an estimated return by \href{https://lilianweng.github.io/posts/2018-02-19-rl-overview/\#monte-carlo-methods}{Monte-Carlo} methods using episode samples to update the policy parameter \$\textbackslash theta\$. REINFORCE works because the expectation of the sample gradient is equal to the actual gradient:

\$\$ \textbackslash begin\{aligned\} \textbackslash nabla\_\textbackslash theta J(\textbackslash theta) \&= \textbackslash mathbb\{E\}\_\textbackslash pi {[}Q\^{}\textbackslash pi(s, a) \textbackslash nabla\_\textbackslash theta \textbackslash ln \textbackslash pi\_\textbackslash theta(a \textbackslash vert s){]} \& \textbackslash\textbackslash{} \&= \textbackslash mathbb\{E\}\_\textbackslash pi {[}G\_t \textbackslash nabla\_\textbackslash theta \textbackslash ln \textbackslash pi\_\textbackslash theta(A\_t \textbackslash vert S\_t){]} \& \textbackslash scriptstyle\{\textbackslash text\{; Because \} Q\^{}\textbackslash pi(S\_t, A\_t) = \textbackslash mathbb\{E\}\_\textbackslash pi{[}G\_t \textbackslash vert S\_t, A\_t{]}\} \textbackslash end\{aligned\} \$\$

Therefore we are able to measure \$G\_t\$ from real sample trajectories and use that to update our policy gradient. It relies on a full trajectory and that's why it is a Monte-Carlo method.

The process is pretty straightforward:

\begin{enumerate}
\tightlist
\item
  Initialize the policy parameter \$\textbackslash theta\$ at random.
\item
  Generate one trajectory on policy \$\textbackslash pi\_\textbackslash theta\$: \$S\_1, A\_1, R\_2, S\_2, A\_2, \textbackslash dots, S\_T\$.
\item
  For t=1, 2, \ldots{} , T:

  \begin{enumerate}
  \tightlist
  \item
    Estimate the the return \$G\_t\$;
  \item
    Update policy parameters: \$\textbackslash theta \textbackslash leftarrow \textbackslash theta + \textbackslash alpha \textbackslash gamma\^{}t G\_t \textbackslash nabla\_\textbackslash theta \textbackslash ln \textbackslash pi\_\textbackslash theta(A\_t \textbackslash vert S\_t)\$
  \end{enumerate}
\end{enumerate}

A widely used variation of REINFORCE is to subtract a baseline value from the return \$G\_t\$ to \emph{reduce the variance of gradient estimation while keeping the bias unchanged} (Remember we always want to do this when possible). For example, a common baseline is to subtract state-value from action-value, and if applied, we would use advantage \$A(s, a) = Q(s, a) - V(s)\$ in the gradient ascent update. This \href{https://danieltakeshi.github.io/2017/03/28/going-deeper-into-reinforcement-learning-fundamentals-of-policy-gradients/}{post} nicely explained why a baseline works for reducing the variance, in addition to a set of fundamentals of policy gradient.

\subsection{\texorpdfstring{Actor-Critic\hyperref[actor-critic]{\#}}{Actor-Critic\#}}\label{actor-critic}

Two main components in policy gradient are the policy model and the value function. It makes a lot of sense to learn the value function in addition to the policy, since knowing the value function can assist the policy update, such as by reducing gradient variance in vanilla policy gradients, and that is exactly what the \textbf{Actor-Critic} method does.

Actor-critic methods consist of two models, which may optionally share parameters:

\begin{itemize}
\tightlist
\item
  \textbf{Critic} updates the value function parameters w and depending on the algorithm it could be action-value \$Q\_w(a \textbackslash vert s)\$ or state-value \$V\_w(s)\$.
\item
  \textbf{Actor} updates the policy parameters \$\textbackslash theta\$ for \$\textbackslash pi\_\textbackslash theta(a \textbackslash vert s)\$, in the direction suggested by the critic.
\end{itemize}

Let's see how it works in a simple action-value actor-critic algorithm.

\begin{enumerate}
\tightlist
\item
  Initialize \$s, \textbackslash theta, w\$ at random; sample \$a \textbackslash sim \textbackslash pi\_\textbackslash theta(a \textbackslash vert s)\$.
\item
  For \$t = 1 \textbackslash dots T\$:

  \begin{enumerate}
  \tightlist
  \item
    Sample reward \$r\_t \textbackslash sim R(s, a)\$ and next state \$s' \textbackslash sim P(s' \textbackslash vert s, a)\$;
  \item
    Then sample the next action \$a' \textbackslash sim \textbackslash pi\_\textbackslash theta(a' \textbackslash vert s')\$;
  \item
    Update the policy parameters: \$\textbackslash theta \textbackslash leftarrow \textbackslash theta + \textbackslash alpha\_\textbackslash theta Q\_w(s, a) \textbackslash nabla\_\textbackslash theta \textbackslash ln \textbackslash pi\_\textbackslash theta(a \textbackslash vert s)\$;
  \item
    Compute the correction (TD error) for action-value at time t:\\
    \$\textbackslash delta\_t = r\_t + \textbackslash gamma Q\_w(s', a') - Q\_w(s, a)\$\\
    and use it to update the parameters of action-value function:\\
    \$w \textbackslash leftarrow w + \textbackslash alpha\_w \textbackslash delta\_t \textbackslash nabla\_w Q\_w(s, a)\$
  \item
    Update \$a \textbackslash leftarrow a'\$ and \$s \textbackslash leftarrow s'\$.
  \end{enumerate}
\end{enumerate}

Two learning rates, \$\textbackslash alpha\_\textbackslash theta\$ and \$\textbackslash alpha\_w\$, are predefined for policy and value function parameter updates respectively.

\subsection{\texorpdfstring{Off-Policy Policy Gradient\hyperref[off-policy-policy-gradient]{\#}}{Off-Policy Policy Gradient\#}}\label{off-policy-policy-gradient}

Both REINFORCE and the vanilla version of actor-critic method are on-policy: training samples are collected according to the target policy --- the very same policy that we try to optimize for. Off policy methods, however, result in several additional advantages:

\begin{enumerate}
\tightlist
\item
  The off-policy approach does not require full trajectories and can reuse any past episodes (\href{https://lilianweng.github.io/posts/2018-02-19-rl-overview/\#deep-q-network}{``experience replay''}) for much better sample efficiency.
\item
  The sample collection follows a behavior policy different from the target policy, bringing better \href{https://lilianweng.github.io/posts/2018-02-19-rl-overview/\#exploration-exploitation-dilemma}{exploration}.
\end{enumerate}

Now let's see how off-policy policy gradient is computed. The behavior policy for collecting samples is a known policy (predefined just like a hyperparameter), labelled as \$\textbackslash beta(a \textbackslash vert s)\$. The objective function sums up the reward over the state distribution defined by this behavior policy:

\$\$ J(\textbackslash theta) = \textbackslash sum\_\{s \textbackslash in \textbackslash mathcal\{S\}\} d\^{}\textbackslash beta(s) \textbackslash sum\_\{a \textbackslash in \textbackslash mathcal\{A\}\} Q\^{}\textbackslash pi(s, a) \textbackslash pi\_\textbackslash theta(a \textbackslash vert s) = \textbackslash mathbb\{E\}\_\{s \textbackslash sim d\^{}\textbackslash beta\} \textbackslash big{[} \textbackslash sum\_\{a \textbackslash in \textbackslash mathcal\{A\}\} Q\^{}\textbackslash pi(s, a) \textbackslash pi\_\textbackslash theta(a \textbackslash vert s) \textbackslash big{]} \$\$

where \$d\^{}\textbackslash beta(s)\$ is the stationary distribution of the behavior policy \$\textbackslash beta\$; recall that \$d\^{}\textbackslash beta(s) = \textbackslash lim\_\{t \textbackslash to \textbackslash infty\} P(S\_t = s \textbackslash vert S\_0, \textbackslash beta)\$; and \$Q\^{}\textbackslash pi\$ is the action-value function estimated with regard to the target policy \$\textbackslash pi\$ (not the behavior policy!).

Given that the training observations are sampled by \$a \textbackslash sim \textbackslash beta(a \textbackslash vert s)\$, we can rewrite the gradient as:

\$\$ \textbackslash begin\{aligned\} \textbackslash nabla\_\textbackslash theta J(\textbackslash theta) \&= \textbackslash nabla\_\textbackslash theta \textbackslash mathbb\{E\}\_\{s \textbackslash sim d\^{}\textbackslash beta\} \textbackslash Big{[} \textbackslash sum\_\{a \textbackslash in \textbackslash mathcal\{A\}\} Q\^{}\textbackslash pi(s, a) \textbackslash pi\_\textbackslash theta(a \textbackslash vert s) \textbackslash Big{]} \& \textbackslash\textbackslash{} \&= \textbackslash mathbb\{E\}\_\{s \textbackslash sim d\^{}\textbackslash beta\} \textbackslash Big{[} \textbackslash sum\_\{a \textbackslash in \textbackslash mathcal\{A\}\} \textbackslash big( Q\^{}\textbackslash pi(s, a) \textbackslash nabla\_\textbackslash theta \textbackslash pi\_\textbackslash theta(a \textbackslash vert s) + \textbackslash color\{red\}\{\textbackslash pi\_\textbackslash theta(a \textbackslash vert s) \textbackslash nabla\_\textbackslash theta Q\^{}\textbackslash pi(s, a)\} \textbackslash big) \textbackslash Big{]} \& \textbackslash scriptstyle\{\textbackslash text\{; Derivative product rule.\}\}\textbackslash\textbackslash{} \&\textbackslash stackrel\{(i)\}\{\textbackslash approx\} \textbackslash mathbb\{E\}\_\{s \textbackslash sim d\^{}\textbackslash beta\} \textbackslash Big{[} \textbackslash sum\_\{a \textbackslash in \textbackslash mathcal\{A\}\} Q\^{}\textbackslash pi(s, a) \textbackslash nabla\_\textbackslash theta \textbackslash pi\_\textbackslash theta(a \textbackslash vert s) \textbackslash Big{]} \& \textbackslash scriptstyle\{\textbackslash text\{; Ignore the red part: \} \textbackslash color\{red\}\{\textbackslash pi\_\textbackslash theta(a \textbackslash vert s) \textbackslash nabla\_\textbackslash theta Q\^{}\textbackslash pi(s, a)\}\}. \textbackslash\textbackslash{} \&= \textbackslash mathbb\{E\}\_\{s \textbackslash sim d\^{}\textbackslash beta\} \textbackslash Big{[} \textbackslash sum\_\{a \textbackslash in \textbackslash mathcal\{A\}\} \textbackslash beta(a \textbackslash vert s) \textbackslash frac\{\textbackslash pi\_\textbackslash theta(a \textbackslash vert s)\}\{\textbackslash beta(a \textbackslash vert s)\} Q\^{}\textbackslash pi(s, a) \textbackslash frac\{\textbackslash nabla\_\textbackslash theta \textbackslash pi\_\textbackslash theta(a \textbackslash vert s)\}\{\textbackslash pi\_\textbackslash theta(a \textbackslash vert s)\} \textbackslash Big{]} \& \textbackslash\textbackslash{} \&= \textbackslash mathbb\{E\}\_\textbackslash beta \textbackslash Big{[}\textbackslash frac\{\textbackslash color\{blue\}\{\textbackslash pi\_\textbackslash theta(a \textbackslash vert s)\}\}\{\textbackslash color\{blue\}\{\textbackslash beta(a \textbackslash vert s)\}\} Q\^{}\textbackslash pi(s, a) \textbackslash nabla\_\textbackslash theta \textbackslash ln \textbackslash pi\_\textbackslash theta(a \textbackslash vert s) \textbackslash Big{]} \& \textbackslash scriptstyle\{\textbackslash text\{; The blue part is the importance weight.\}\} \textbackslash end\{aligned\} \$\$

where \$\textbackslash frac\{\textbackslash pi\_\textbackslash theta(a \textbackslash vert s)\}\{\textbackslash beta(a \textbackslash vert s)\}\$ is the \href{http://timvieira.github.io/blog/post/2014/12/21/importance-sampling/}{importance weight}. Because \$Q\^{}\textbackslash pi\$ is a function of the target policy and thus a function of policy parameter \$\textbackslash theta\$, we should take the derivative of \$\textbackslash nabla\_\textbackslash theta Q\^{}\textbackslash pi(s, a)\$ as well according to the product rule. However, it is super hard to compute \$\textbackslash nabla\_\textbackslash theta Q\^{}\textbackslash pi(s, a)\$ in reality. Fortunately if we use an approximated gradient with the gradient of Q ignored, we still guarantee the policy improvement and eventually achieve the true local minimum. This is justified in the proof \href{https://arxiv.org/pdf/1205.4839.pdf}{here} (Degris, White \& Sutton, 2012).

In summary, when applying policy gradient in the off-policy setting, we can simple adjust it with a weighted sum and the weight is the ratio of the target policy to the behavior policy, \$\textbackslash frac\{\textbackslash pi\_\textbackslash theta(a \textbackslash vert s)\}\{\textbackslash beta(a \textbackslash vert s)\}\$.

\subsection{\texorpdfstring{A3C\hyperref[a3c]{\#}}{A3C\#}}\label{a3c}

{[}\href{https://arxiv.org/abs/1602.01783}{paper}\textbar{}\href{https://github.com/dennybritz/reinforcement-learning/tree/master/PolicyGradient/a3c}{code}{]}

\textbf{Asynchronous Advantage Actor-Critic} (\href{https://arxiv.org/abs/1602.01783}{Mnih et al., 2016}), short for \textbf{A3C}, is a classic policy gradient method with a special focus on parallel training.

In A3C, the critics learn the value function while multiple actors are trained in parallel and get synced with global parameters from time to time. Hence, A3C is designed to work well for parallel training.

Let's use the state-value function as an example. The loss function for state value is to minimize the mean squared error, \$J\_v(w) = (G\_t - V\_w(s))\^{}2\$ and gradient descent can be applied to find the optimal w. This state-value function is used as the baseline in the policy gradient update.

Here is the algorithm outline:

\begin{enumerate}
\item
  We have global parameters, \$\textbackslash theta\$ and \$w\$; similar thread-specific parameters, \$\textbackslash theta'\$ and \$w'\$.
\item
  Initialize the time step \$t = 1\$
\item
  While \$T \textbackslash leq T\_\textbackslash text\{MAX\}\$:

  \begin{enumerate}
  \tightlist
  \item
    Reset gradient: \$\textbackslash mathrm\{d\}\textbackslash theta = 0\$ and \$\textbackslash mathrm\{d\}w = 0\$.
  \item
    Synchronize thread-specific parameters with global ones: \$\textbackslash theta' = \textbackslash theta\$ and \$w' = w\$.
  \item
    \$t\_\textbackslash text\{start\}\$ = t and sample a starting state \$s\_t\$.
  \item
    While (\$s\_t\$ != TERMINAL) and \$t - t\_\textbackslash text\{start\} \textbackslash leq t\_\textbackslash text\{max\}\$:

    \begin{enumerate}
    \tightlist
    \item
      Pick the action \$A\_t \textbackslash sim \textbackslash pi\_\{\textbackslash theta'\}(A\_t \textbackslash vert S\_t)\$ and receive a new reward \$R\_t\$ and a new state \$s\_\{t+1\}\$.
    \item
      Update \$t = t + 1\$ and \$T = T + 1\$
    \end{enumerate}
  \item
    Initialize the variable that holds the return estimation
  \end{enumerate}

  \$\$ R = \textbackslash begin\{cases\} 0 \& \textbackslash text\{if \} s\_t \textbackslash text\{ is TERMINAL\} \textbackslash\textbackslash{} V\_\{w\textquotesingle\}(s\_t) \& \textbackslash text\{otherwise\} \textbackslash end\{cases\} \$\$

  6. For \$i = t-1, \textbackslash dots, t\textbackslash\_\textbackslash text\{start\}\$: 1. \$R \textbackslash leftarrow \textbackslash gamma R + R\textbackslash\_i\$; here R is a MC measure of \$G\textbackslash\_i\$. 2. Accumulate gradients w.r.t. \$\textbackslash theta\textquotesingle\$: \$d\textbackslash theta \textbackslash leftarrow d\textbackslash theta + \textbackslash nabla\textbackslash\_\{\textbackslash theta\textquotesingle\} \textbackslash log \textbackslash pi\textbackslash\_\{\textbackslash theta\textquotesingle\}(a\textbackslash\_i \textbackslash vert s\textbackslash\_i)(R - V\textbackslash\_\{w\textquotesingle\}(s\textbackslash\_i))\$;\\
  Accumulate gradients w.r.t. w\textquotesingle: \$dw \textbackslash leftarrow dw + 2 (R - V\textbackslash\_\{w\textquotesingle\}(s\textbackslash\_i)) \textbackslash nabla\textbackslash\_\{w\textquotesingle\} (R - V\textbackslash\_\{w\textquotesingle\}(s\textbackslash\_i))\$.

  \begin{enumerate}
  \setcounter{enumii}{6}
  \tightlist
  \item
    Update asynchronously \$\textbackslash theta\$ using \$\textbackslash mathrm\{d\}\textbackslash theta\$, and \$w\$ using \$\textbackslash mathrm\{d\}w\$.
  \end{enumerate}
\end{enumerate}

A3C enables the parallelism in multiple agent training. The gradient accumulation step (6.2) can be considered as a parallelized reformation of minibatch-based stochastic gradient update: the values of \$w\$ or \$\textbackslash theta\$ get corrected by a little bit in the direction of each training thread independently.

\subsection{\texorpdfstring{A2C\hyperref[a2c]{\#}}{A2C\#}}\label{a2c}

{[}\href{https://arxiv.org/abs/1602.01783}{paper}\textbar{}\href{https://github.com/openai/baselines/blob/master/baselines/a2c/a2c.py}{code}{]}

\textbf{A2C} is a synchronous, deterministic version of A3C; that's why it is named as ``A2C'' with the first ``A'' (``asynchronous'') removed. In A3C each agent talks to the global parameters independently, so it is possible sometimes the thread-specific agents would be playing with policies of different versions and therefore the aggregated update would not be optimal. To resolve the inconsistency, a coordinator in A2C waits for all the parallel actors to finish their work before updating the global parameters and then in the next iteration parallel actors starts from the same policy. The synchronized gradient update keeps the training more cohesive and potentially to make convergence faster.

A2C has been \href{https://blog.openai.com/baselines-acktr-a2c/}{shown} to be able to utilize GPUs more efficiently and work better with large batch sizes while achieving same or better performance than A3C.

\begin{figure}
\centering
\includegraphics{A3C_vs_A2C.png}
\caption{The architecture of A3C versus A2C.}
\end{figure}

\subsection{\texorpdfstring{DPG\hyperref[dpg]{\#}}{DPG\#}}\label{dpg}

{[}\href{https://hal.inria.fr/file/index/docid/938992/filename/dpg-icml2014.pdf}{paper}\textbar code{]}

In methods described above, the policy function \$\textbackslash pi(. \textbackslash vert s)\$ is always modeled as a probability distribution over actions \$\textbackslash mathcal\{A\}\$ given the current state and thus it is \emph{stochastic}. \textbf{Deterministic policy gradient (DPG)} instead models the policy as a deterministic decision: \$a = \textbackslash mu(s)\$. It may look bizarre --- how can you calculate the gradient of the action probability when it outputs a single action? Let's look into it step by step.

Refresh on a few notations to facilitate the discussion:

\begin{itemize}
\tightlist
\item
  \$\textbackslash rho\_0(s)\$: The initial distribution over states
\item
  \$\textbackslash rho\^{}\textbackslash mu(s \textbackslash to s', k)\$: Starting from state s, the visitation probability density at state s' after moving k steps by policy \$\textbackslash mu\$.
\item
  \$\textbackslash rho\^{}\textbackslash mu(s')\$: Discounted state distribution, defined as \$\textbackslash rho\^{}\textbackslash mu(s') = \textbackslash int\_\textbackslash mathcal\{S\} \textbackslash sum\_\{k=1\}\^{}\textbackslash infty \textbackslash gamma\^{}\{k-1\} \textbackslash rho\_0(s) \textbackslash rho\^{}\textbackslash mu(s \textbackslash to s', k) ds\$.
\end{itemize}

The objective function to optimize for is listed as follows:

\$\$ J(\textbackslash theta) = \textbackslash int\_\textbackslash mathcal\{S\} \textbackslash rho\^{}\textbackslash mu(s) Q(s, \textbackslash mu\_\textbackslash theta(s)) ds \$\$

\textbf{Deterministic policy gradient theorem}: Now it is the time to compute the gradient! According to the chain rule, we first take the gradient of Q w.r.t. the action a and then take the gradient of the deterministic policy function \$\textbackslash mu\$ w.r.t. \$\textbackslash theta\$:

\$\$ \textbackslash begin\{aligned\} \textbackslash nabla\_\textbackslash theta J(\textbackslash theta) \&= \textbackslash int\_\textbackslash mathcal\{S\} \textbackslash rho\^{}\textbackslash mu(s) \textbackslash nabla\_a Q\^{}\textbackslash mu(s, a) \textbackslash nabla\_\textbackslash theta \textbackslash mu\_\textbackslash theta(s) \textbackslash rvert\_\{a=\textbackslash mu\_\textbackslash theta(s)\} ds \textbackslash\textbackslash{} \&= \textbackslash mathbb\{E\}\_\{s \textbackslash sim \textbackslash rho\^{}\textbackslash mu\} {[}\textbackslash nabla\_a Q\^{}\textbackslash mu(s, a) \textbackslash nabla\_\textbackslash theta \textbackslash mu\_\textbackslash theta(s) \textbackslash rvert\_\{a=\textbackslash mu\_\textbackslash theta(s)\}{]} \textbackslash end\{aligned\} \$\$

We can consider the deterministic policy as a \emph{special case} of the stochastic one, when the probability distribution contains only one extreme non-zero value over one action. Actually, in the DPG \href{https://hal.inria.fr/file/index/docid/938992/filename/dpg-icml2014.pdf}{paper}, the authors have shown that if the stochastic policy \$\textbackslash pi\_\{\textbackslash mu\_\textbackslash theta, \textbackslash sigma\}\$ is re-parameterized by a deterministic policy \$\textbackslash mu\_\textbackslash theta\$ and a variation variable \$\textbackslash sigma\$, the stochastic policy is eventually equivalent to the deterministic case when \$\textbackslash sigma=0\$. Compared to the deterministic policy, we expect the stochastic policy to require more samples as it integrates the data over the whole state and action space.

The deterministic policy gradient theorem can be plugged into common policy gradient frameworks.

Let's consider an example of on-policy actor-critic algorithm to showcase the procedure. In each iteration of on-policy actor-critic, two actions are taken deterministically \$a = \textbackslash mu\_\textbackslash theta(s)\$ and the \href{https://lilianweng.github.io/posts/2018-02-19-rl-overview/\#sarsa-on-policy-td-control}{SARSA} update on policy parameters relies on the new gradient that we just computed above:

\$\$ \textbackslash begin\{aligned\} \textbackslash delta\_t \&= R\_t + \textbackslash gamma Q\_w(s\_\{t+1\}, a\_\{t+1\}) - Q\_w(s\_t, a\_t) \& \textbackslash small\{\textbackslash text\{; TD error in SARSA\}\}\textbackslash\textbackslash{} w\_\{t+1\} \&= w\_t + \textbackslash alpha\_w \textbackslash delta\_t \textbackslash nabla\_w Q\_w(s\_t, a\_t) \& \textbackslash\textbackslash{} \textbackslash theta\_\{t+1\} \&= \textbackslash theta\_t + \textbackslash alpha\_\textbackslash theta \textbackslash color\{red\}\{\textbackslash nabla\_a Q\_w(s\_t, a\_t) \textbackslash nabla\_\textbackslash theta \textbackslash mu\_\textbackslash theta(s) \textbackslash rvert\_\{a=\textbackslash mu\_\textbackslash theta(s)\}\} \& \textbackslash small\{\textbackslash text\{; Deterministic policy gradient theorem\}\} \textbackslash end\{aligned\} \$\$

However, unless there is sufficient noise in the environment, it is very hard to guarantee enough \href{https://lilianweng.github.io/posts/2018-02-19-rl-overview/\#exploration-exploitation-dilemma}{exploration} due to the determinacy of the policy. We can either add noise into the policy (ironically this makes it nondeterministic!) or learn it off-policy-ly by following a different stochastic behavior policy to collect samples.

Say, in the off-policy approach, the training trajectories are generated by a stochastic policy \$\textbackslash beta(a \textbackslash vert s)\$ and thus the state distribution follows the corresponding discounted state density \$\textbackslash rho\^{}\textbackslash beta\$:

\$\$ \textbackslash begin\{aligned\} J\_\textbackslash beta(\textbackslash theta) \&= \textbackslash int\_\textbackslash mathcal\{S\} \textbackslash rho\^{}\textbackslash beta Q\^{}\textbackslash mu(s, \textbackslash mu\_\textbackslash theta(s)) ds \textbackslash\textbackslash{} \textbackslash nabla\_\textbackslash theta J\_\textbackslash beta(\textbackslash theta) \&= \textbackslash mathbb\{E\}\_\{s \textbackslash sim \textbackslash rho\^{}\textbackslash beta\} {[}\textbackslash nabla\_a Q\^{}\textbackslash mu(s, a) \textbackslash nabla\_\textbackslash theta \textbackslash mu\_\textbackslash theta(s) \textbackslash rvert\_\{a=\textbackslash mu\_\textbackslash theta(s)\} {]} \textbackslash end\{aligned\} \$\$

Note that because the policy is deterministic, we only need \$Q\^{}\textbackslash mu(s, \textbackslash mu\_\textbackslash theta(s))\$ rather than \$\textbackslash sum\_a \textbackslash pi(a \textbackslash vert s) Q\^{}\textbackslash pi(s, a)\$ as the estimated reward of a given state s. In the off-policy approach with a stochastic policy, importance sampling is often used to correct the mismatch between behavior and target policies, as what we have described \hyperref[off-policy-policy-gradient]{above}. However, because the deterministic policy gradient removes the integral over actions, we can avoid importance sampling.

\subsection{\texorpdfstring{DDPG\hyperref[ddpg]{\#}}{DDPG\#}}\label{ddpg}

{[}\href{https://arxiv.org/pdf/1509.02971.pdf}{paper}\textbar{}\href{https://github.com/openai/baselines/tree/master/baselines/ddpg}{code}{]}

\textbf{DDPG} (\href{https://arxiv.org/pdf/1509.02971.pdf}{Lillicrap, et al., 2015}), short for \textbf{Deep Deterministic Policy Gradient}, is a model-free off-policy actor-critic algorithm, combining \hyperref[dpg]{DPG} with \href{https://lilianweng.github.io/posts/2018-02-19-rl-overview/\#deep-q-network}{DQN}. Recall that DQN (Deep Q-Network) stabilizes the learning of Q-function by experience replay and the frozen target network. The original DQN works in discrete space, and DDPG extends it to continuous space with the actor-critic framework while learning a deterministic policy.

In order to do better exploration, an exploration policy \$\textbackslash mu'\$ is constructed by adding noise \$\textbackslash mathcal\{N\}\$:

\$\$ \textbackslash mu\textquotesingle(s) = \textbackslash mu\_\textbackslash theta(s) + \textbackslash mathcal\{N\} \$\$

In addition, DDPG does soft updates (``conservative policy iteration'') on the parameters of both actor and critic, with \$\textbackslash tau \textbackslash ll 1\$: \$\textbackslash theta' \textbackslash leftarrow \textbackslash tau \textbackslash theta + (1 - \textbackslash tau) \textbackslash theta'\$. In this way, the target network values are constrained to change slowly, different from the design in DQN that the target network stays frozen for some period of time.

One detail in the paper that is particularly useful in robotics is on how to normalize the different physical units of low dimensional features. For example, a model is designed to learn a policy with the robot's positions and velocities as input; these physical statistics are different by nature and even statistics of the same type may vary a lot across multiple robots. \href{http://proceedings.mlr.press/v37/ioffe15.pdf}{Batch normalization} is applied to fix it by normalizing every dimension across samples in one minibatch.

\begin{figure}
\centering
\includegraphics{DDPG_algo.png}
\caption{Fig 3. DDPG Algorithm. (Image source: \href{https://arxiv.org/pdf/1509.02971.pdf}{Lillicrap, et al., 2015})}
\end{figure}

\subsection{\texorpdfstring{D4PG\hyperref[d4pg]{\#}}{D4PG\#}}\label{d4pg}

{[}\href{https://openreview.net/forum?id=SyZipzbCb}{paper}\textbar code (Search ``github d4pg'' and you will see a few.){]}

\textbf{Distributed Distributional DDPG (D4PG)} applies a set of improvements on DDPG to make it run in the distributional fashion.

(1) \textbf{Distributional Critic}: The critic estimates the expected Q value as a random variable \textasciitilde{} a distribution \$Z\_w\$ parameterized by \$w\$ and therefore \$Q\_w(s, a) = \textbackslash mathbb\{E\} Z\_w(x, a)\$. The loss for learning the distribution parameter is to minimize some measure of the distance between two distributions --- distributional TD error: \$L(w) = \textbackslash mathbb\{E\}{[}d(\textbackslash mathcal\{T\}\_\{\textbackslash mu\_\textbackslash theta\}, Z\_\{w'\}(s, a), Z\_w(s, a){]}\$, where \$\textbackslash mathcal\{T\}\_\{\textbackslash mu\_\textbackslash theta\}\$ is the Bellman operator.

The deterministic policy gradient update becomes:

\$\$ \textbackslash begin\{aligned\} \textbackslash nabla\_\textbackslash theta J(\textbackslash theta) \&\textbackslash approx \textbackslash mathbb\{E\}\_\{\textbackslash rho\^{}\textbackslash mu\} {[}\textbackslash nabla\_a Q\_w(s, a) \textbackslash nabla\_\textbackslash theta \textbackslash mu\_\textbackslash theta(s) \textbackslash rvert\_\{a=\textbackslash mu\_\textbackslash theta(s)\}{]} \& \textbackslash scriptstyle\{\textbackslash text\{; gradient update in DPG\}\} \textbackslash\textbackslash{} \&= \textbackslash mathbb\{E\}\_\{\textbackslash rho\^{}\textbackslash mu\} {[}\textbackslash mathbb\{E\}{[}\textbackslash nabla\_a Z\_w(s, a){]} \textbackslash nabla\_\textbackslash theta \textbackslash mu\_\textbackslash theta(s) \textbackslash rvert\_\{a=\textbackslash mu\_\textbackslash theta(s)\}{]} \& \textbackslash scriptstyle\{\textbackslash text\{; expectation of the Q-value distribution.\}\} \textbackslash end\{aligned\} \$\$

(2) \textbf{\$N\$-step returns}: When calculating the TD error, D4PG computes \$N\$-step TD target rather than one-step to incorporate rewards in more future steps. Thus the new TD target is:

\$\$ r(s\_0, a\_0) + \textbackslash mathbb\{E\}{[}\textbackslash sum\_\{n=1\}\^{}\{N-1\} r(s\_n, a\_n) + \textbackslash gamma\^{}N Q(s\_N, \textbackslash mu\_\textbackslash theta(s\_N)) \textbackslash vert s\_0, a\_0 {]} \$\$

(3) \textbf{Multiple Distributed Parallel Actors}: D4PG utilizes \$K\$ independent actors, gathering experience in parallel and feeding data into the same replay buffer.

(4) \textbf{Prioritized Experience Replay (\href{https://arxiv.org/abs/1511.05952}{PER})}: The last piece of modification is to do sampling from the replay buffer of size \$R\$ with an non-uniform probability \$p\_i\$. In this way, a sample \$i\$ has the probability \$(Rp\_i)\^{}\{-1\}\$ to be selected and thus the importance weight is \$(Rp\_i)\^{}\{-1\}\$.

\begin{figure}
\centering
\includegraphics{D4PG_algo.png}
\caption{D4PG algorithm (Image source: \href{https://openreview.net/forum?id=SyZipzbCb}{Barth-Maron, et al. 2018}); Note that in the original paper, the variable letters are chosen slightly differently from what in the post; i.e. I use \$\textbackslash mu(.)\$ for representing a deterministic policy instead of \$\textbackslash pi(.)\$.}
\end{figure}

\subsection{\texorpdfstring{MADDPG\hyperref[maddpg]{\#}}{MADDPG\#}}\label{maddpg}

{[}\href{https://arxiv.org/pdf/1706.02275.pdf}{paper}\textbar{}\href{https://github.com/openai/maddpg}{code}{]}

\textbf{Multi-agent DDPG} (\textbf{MADDPG}) (\href{https://arxiv.org/pdf/1706.02275.pdf}{Lowe et al., 2017}) extends DDPG to an environment where multiple agents are coordinating to complete tasks with only local information. In the viewpoint of one agent, the environment is non-stationary as policies of other agents are quickly upgraded and remain unknown. MADDPG is an actor-critic model redesigned particularly for handling such a changing environment and interactions between agents.

The problem can be formalized in the multi-agent version of MDP, also known as \emph{Markov games}. MADDPG is proposed for partially observable Markov games. Say, there are N agents in total with a set of states \$\textbackslash mathcal\{S\}\$. Each agent owns a set of possible action, \$\textbackslash mathcal\{A\}\_1, \textbackslash dots, \textbackslash mathcal\{A\}\_N\$, and a set of observation, \$\textbackslash mathcal\{O\}\_1, \textbackslash dots, \textbackslash mathcal\{O\}\_N\$. The state transition function involves all states, action and observation spaces \$\textbackslash mathcal\{T\}: \textbackslash mathcal\{S\} \textbackslash times \textbackslash mathcal\{A\}\_1 \textbackslash times \textbackslash dots \textbackslash mathcal\{A\}\_N \textbackslash mapsto \textbackslash mathcal\{S\}\$. Each agent's stochastic policy only involves its own state and action: \$\textbackslash pi\_\{\textbackslash theta\_i\}: \textbackslash mathcal\{O\}\_i \textbackslash times \textbackslash mathcal\{A\}\_i \textbackslash mapsto {[}0, 1{]}\$, a probability distribution over actions given its own observation, or a deterministic policy: \$\textbackslash mu\_\{\textbackslash theta\_i\}: \textbackslash mathcal\{O\}\_i \textbackslash mapsto \textbackslash mathcal\{A\}\_i\$.

Let \$\textbackslash vec\{o\} = \{o\_1, \textbackslash dots, o\_N\}\$, \$\textbackslash vec\{\textbackslash mu\} = \{\textbackslash mu\_1, \textbackslash dots, \textbackslash mu\_N\}\$ and the policies are parameterized by \$\textbackslash vec\{\textbackslash theta\} = \{\textbackslash theta\_1, \textbackslash dots, \textbackslash theta\_N\}\$.

The critic in MADDPG learns a centralized action-value function \$Q\^{}\textbackslash vec\{\textbackslash mu\}\_i(\textbackslash vec\{o\}, a\_1, \textbackslash dots, a\_N)\$ for the i-th agent, where \$a\_1 \textbackslash in \textbackslash mathcal\{A\}\_1, \textbackslash dots, a\_N \textbackslash in \textbackslash mathcal\{A\}\_N\$ are actions of all agents. Each \$Q\^{}\textbackslash vec\{\textbackslash mu\}\_i\$ is learned separately for \$i=1, \textbackslash dots, N\$ and therefore multiple agents can have arbitrary reward structures, including conflicting rewards in a competitive setting. Meanwhile, multiple actors, one for each agent, are exploring and upgrading the policy parameters \$\textbackslash theta\_i\$ on their own.

\textbf{Actor update}:

\$\$ \textbackslash nabla\_\{\textbackslash theta\_i\} J(\textbackslash theta\_i) = \textbackslash mathbb\{E\}\_\{\textbackslash vec\{o\}, a \textbackslash sim \textbackslash mathcal\{D\}\} {[}\textbackslash nabla\_\{a\_i\} Q\^{}\{\textbackslash vec\{\textbackslash mu\}\}\_i (\textbackslash vec\{o\}, a\_1, \textbackslash dots, a\_N) \textbackslash nabla\_\{\textbackslash theta\_i\} \textbackslash mu\_\{\textbackslash theta\_i\}(o\_i) \textbackslash rvert\_\{a\_i=\textbackslash mu\_\{\textbackslash theta\_i\}(o\_i)\} {]} \$\$

Where \$\textbackslash mathcal\{D\}\$ is the memory buffer for experience replay, containing multiple episode samples \$(\textbackslash vec\{o\}, a\_1, \textbackslash dots, a\_N, r\_1, \textbackslash dots, r\_N, \textbackslash vec\{o\}')\$ --- given current observation \$\textbackslash vec\{o\}\$, agents take action \$a\_1, \textbackslash dots, a\_N\$ and get rewards \$r\_1, \textbackslash dots, r\_N\$, leading to the new observation \$\textbackslash vec\{o\}'\$.

\textbf{Critic update}:

\$\$ \textbackslash begin\{aligned\} \textbackslash mathcal\{L\}(\textbackslash theta\_i) \&= \textbackslash mathbb\{E\}\_\{\textbackslash vec\{o\}, a\_1, \textbackslash dots, a\_N, r\_1, \textbackslash dots, r\_N, \textbackslash vec\{o\}\textquotesingle\}{[} (Q\^{}\{\textbackslash vec\{\textbackslash mu\}\}\_i(\textbackslash vec\{o\}, a\_1, \textbackslash dots, a\_N) - y)\^{}2 {]} \& \textbackslash\textbackslash{} \textbackslash text\{where \} y \&= r\_i + \textbackslash gamma Q\^{}\{\textbackslash vec\{\textbackslash mu\}\textquotesingle\}\_i (\textbackslash vec\{o\}\textquotesingle, a\textquotesingle\_1, \textbackslash dots, a\textquotesingle\_N) \textbackslash rvert\_\{a\textquotesingle\_j = \textbackslash mu\textquotesingle\_\{\textbackslash theta\_j\}\} \& \textbackslash scriptstyle\{\textbackslash text\{; TD target!\}\} \textbackslash end\{aligned\} \$\$

where \$\textbackslash vec\{\textbackslash mu\}'\$ are the target policies with delayed softly-updated parameters.

If the policies \$\textbackslash vec\{\textbackslash mu\}\$ are unknown during the critic update, we can ask each agent to learn and evolve its own approximation of others' policies. Using the approximated policies, MADDPG still can learn efficiently although the inferred policies might not be accurate.

To mitigate the high variance triggered by the interaction between competing or collaborating agents in the environment, MADDPG proposed one more element - \emph{policy ensembles}:

\begin{enumerate}
\tightlist
\item
  Train K policies for one agent;
\item
  Pick a random policy for episode rollouts;
\item
  Take an ensemble of these K policies to do gradient update.
\end{enumerate}

In summary, MADDPG added three additional ingredients on top of DDPG to make it adapt to the multi-agent environment:

\begin{itemize}
\tightlist
\item
  Centralized critic + decentralized actors;
\item
  Actors are able to use estimated policies of other agents for learning;
\item
  Policy ensembling is good for reducing variance.
\end{itemize}

\begin{figure}
\centering
\includegraphics{MADDPG.png}
\caption{The architecture design of MADDPG. (Image source: \href{https://arxiv.org/pdf/1706.02275.pdf}{Lowe et al., 2017})}
\end{figure}

\subsection{\texorpdfstring{TRPO\hyperref[trpo]{\#}}{TRPO\#}}\label{trpo}

{[}\href{https://arxiv.org/pdf/1502.05477.pdf}{paper}\textbar{}\href{https://github.com/openai/baselines/tree/master/baselines/trpo_mpi}{code}{]}

To improve training stability, we should avoid parameter updates that change the policy too much at one step. \textbf{Trust region policy optimization (TRPO)} (\href{https://arxiv.org/pdf/1502.05477.pdf}{Schulman, et al., 2015}) carries out this idea by enforcing a \href{https://lilianweng.github.io/posts/2017-08-20-gan/\#kullbackleibler-and-jensenshannon-divergence}{KL divergence} constraint on the size of policy update at each iteration.

Consider the case when we are doing off-policy RL, the policy \$\textbackslash beta\$ used for collecting trajectories on rollout workers is different from the policy \$\textbackslash pi\$ to optimize for. The objective function in an off-policy model measures the total advantage over the state visitation distribution and actions, while the mismatch between the training data distribution and the true policy state distribution is compensated by importance sampling estimator:

\$\$ \textbackslash begin\{aligned\} J(\textbackslash theta) \&= \textbackslash sum\_\{s \textbackslash in \textbackslash mathcal\{S\}\} \textbackslash rho\^{}\{\textbackslash pi\_\{\textbackslash theta\_\textbackslash text\{old\}\}\} \textbackslash sum\_\{a \textbackslash in \textbackslash mathcal\{A\}\} \textbackslash big( \textbackslash pi\_\textbackslash theta(a \textbackslash vert s) \textbackslash hat\{A\}\_\{\textbackslash theta\_\textbackslash text\{old\}\}(s, a) \textbackslash big) \& \textbackslash\textbackslash{} \&= \textbackslash sum\_\{s \textbackslash in \textbackslash mathcal\{S\}\} \textbackslash rho\^{}\{\textbackslash pi\_\{\textbackslash theta\_\textbackslash text\{old\}\}\} \textbackslash sum\_\{a \textbackslash in \textbackslash mathcal\{A\}\} \textbackslash big( \textbackslash beta(a \textbackslash vert s) \textbackslash frac\{\textbackslash pi\_\textbackslash theta(a \textbackslash vert s)\}\{\textbackslash beta(a \textbackslash vert s)\} \textbackslash hat\{A\}\_\{\textbackslash theta\_\textbackslash text\{old\}\}(s, a) \textbackslash big) \& \textbackslash scriptstyle\{\textbackslash text\{; Importance sampling\}\} \textbackslash\textbackslash{} \&= \textbackslash mathbb\{E\}\_\{s \textbackslash sim \textbackslash rho\^{}\{\textbackslash pi\_\{\textbackslash theta\_\textbackslash text\{old\}\}\}, a \textbackslash sim \textbackslash beta\} \textbackslash big{[} \textbackslash frac\{\textbackslash pi\_\textbackslash theta(a \textbackslash vert s)\}\{\textbackslash beta(a \textbackslash vert s)\} \textbackslash hat\{A\}\_\{\textbackslash theta\_\textbackslash text\{old\}\}(s, a) \textbackslash big{]} \& \textbackslash end\{aligned\} \$\$

where \$\textbackslash theta\_\textbackslash text\{old\}\$ is the policy parameters before the update and thus known to us; \$\textbackslash rho\^{}\{\textbackslash pi\_\{\textbackslash theta\_\textbackslash text\{old\}\}\}\$ is defined in the same way as \hyperref[dpg]{above}; \$\textbackslash beta(a \textbackslash vert s)\$ is the behavior policy for collecting trajectories. Noted that we use an estimated advantage \$\textbackslash hat\{A\}(.)\$ rather than the true advantage function \$A(.)\$ because the true rewards are usually unknown.

When training on policy, theoretically the policy for collecting data is same as the policy that we want to optimize. However, when rollout workers and optimizers are running in parallel asynchronously, the behavior policy can get stale. TRPO considers this subtle difference: It labels the behavior policy as \$\textbackslash pi\_\{\textbackslash theta\_\textbackslash text\{old\}\}(a \textbackslash vert s)\$ and thus the objective function becomes:

\$\$ J(\textbackslash theta) = \textbackslash mathbb\{E\}\_\{s \textbackslash sim \textbackslash rho\^{}\{\textbackslash pi\_\{\textbackslash theta\_\textbackslash text\{old\}\}\}, a \textbackslash sim \textbackslash pi\_\{\textbackslash theta\_\textbackslash text\{old\}\}\} \textbackslash big{[} \textbackslash frac\{\textbackslash pi\_\textbackslash theta(a \textbackslash vert s)\}\{\textbackslash pi\_\{\textbackslash theta\_\textbackslash text\{old\}\}(a \textbackslash vert s)\} \textbackslash hat\{A\}\_\{\textbackslash theta\_\textbackslash text\{old\}\}(s, a) \textbackslash big{]} \$\$

TRPO aims to maximize the objective function \$J(\textbackslash theta)\$ subject to, \emph{trust region constraint} which enforces the distance between old and new policies measured by \href{https://en.wikipedia.org/wiki/Kullback\%E2\%80\%93Leibler_divergence}{KL-divergence} to be small enough, within a parameter δ:

\$\$ \textbackslash mathbb\{E\}\_\{s \textbackslash sim \textbackslash rho\^{}\{\textbackslash pi\_\{\textbackslash theta\_\textbackslash text\{old\}\}\}\} {[}D\_\textbackslash text\{KL\}(\textbackslash pi\_\{\textbackslash theta\_\textbackslash text\{old\}\}(.\textbackslash vert s) \textbackslash\textbar{} \textbackslash pi\_\textbackslash theta(.\textbackslash vert s){]} \textbackslash leq \textbackslash delta \$\$

In this way, the old and new policies would not diverge too much when this hard constraint is met. While still, TRPO can guarantee a monotonic improvement over policy iteration (Neat, right?). Please read the proof in the \href{https://arxiv.org/pdf/1502.05477.pdf}{paper} if interested :)

\subsection{\texorpdfstring{PPO\hyperref[ppo]{\#}}{PPO\#}}\label{ppo}

{[}\href{https://arxiv.org/pdf/1707.06347.pdf}{paper}\textbar{}\href{https://github.com/openai/baselines/tree/master/baselines/ppo1}{code}{]}

Given that TRPO is relatively complicated and we still want to implement a similar constraint, \textbf{proximal policy optimization (PPO)} simplifies it by using a clipped surrogate objective while retaining similar performance.

First, let's denote the probability ratio between old and new policies as:

\$\$ r(\textbackslash theta) = \textbackslash frac\{\textbackslash pi\_\textbackslash theta(a \textbackslash vert s)\}\{\textbackslash pi\_\{\textbackslash theta\_\textbackslash text\{old\}\}(a \textbackslash vert s)\} \$\$

Then, the objective function of TRPO (on policy) becomes:

\$\$ J\^{}\textbackslash text\{TRPO\} (\textbackslash theta) = \textbackslash mathbb\{E\} {[} r(\textbackslash theta) \textbackslash hat\{A\}\_\{\textbackslash theta\_\textbackslash text\{old\}\}(s, a) {]} \$\$

\phantomsection\label{ppo_loss}{}Without a limitation on the distance between \$\textbackslash theta\_\textbackslash text\{old\}\$ and \$\textbackslash theta\$, to maximize \$J\^{}\textbackslash text\{TRPO\} (\textbackslash theta)\$ would lead to instability with extremely large parameter updates and big policy ratios. PPO imposes the constraint by forcing \$r(\textbackslash theta)\$ to stay within a small interval around 1, precisely \${[}1-\textbackslash epsilon, 1+\textbackslash epsilon{]}\$, where \$\textbackslash epsilon\$ is a hyperparameter.

\$\$ J\^{}\textbackslash text\{CLIP\} (\textbackslash theta) = \textbackslash mathbb\{E\} {[} \textbackslash min( r(\textbackslash theta) \textbackslash hat\{A\}\_\{\textbackslash theta\_\textbackslash text\{old\}\}(s, a), \textbackslash text\{clip\}(r(\textbackslash theta), 1 - \textbackslash epsilon, 1 + \textbackslash epsilon) \textbackslash hat\{A\}\_\{\textbackslash theta\_\textbackslash text\{old\}\}(s, a)){]} \$\$

The function \$\textbackslash text\{clip\}(r(\textbackslash theta), 1 - \textbackslash epsilon, 1 + \textbackslash epsilon)\$ clips the ratio to be no more than \$1+\textbackslash epsilon\$ and no less than \$1-\textbackslash epsilon\$. The objective function of PPO takes the minimum one between the original value and the clipped version and therefore we lose the motivation for increasing the policy update to extremes for better rewards.

When applying PPO on the network architecture with shared parameters for both policy (actor) and value (critic) functions, in addition to the clipped reward, the objective function is augmented with an error term on the value estimation (formula in red) and an entropy term (formula in blue) to encourage sufficient exploration.

\$\$ J\^{}\textbackslash text\{CLIP\textquotesingle\} (\textbackslash theta) = \textbackslash mathbb\{E\} {[} J\^{}\textbackslash text\{CLIP\} (\textbackslash theta) - \textbackslash color\{red\}\{c\_1 (V\_\textbackslash theta(s) - V\_\textbackslash text\{target\})\^{}2\} + \textbackslash color\{blue\}\{c\_2 H(s, \textbackslash pi\_\textbackslash theta(.))\} {]} \$\$

where Both \$c\_1\$ and \$c\_2\$ are two hyperparameter constants.

PPO has been tested on a set of benchmark tasks and proved to produce awesome results with much greater simplicity.

In a later paper by \href{https://arxiv.org/abs/2009.10897}{Hsu et al., 2020}, two common design choices in PPO are revisited, precisely (1) clipped probability ratio for policy regularization and (2) parameterize policy action space by continuous Gaussian or discrete softmax distribution. They first identified three failure modes in PPO and proposed replacements for these two designs.

The failure modes are:

\begin{enumerate}
\tightlist
\item
  On continuous action spaces, standard PPO is unstable when rewards vanish outside bounded support.
\item
  On discrete action spaces with sparse high rewards, standard PPO often gets stuck at suboptimal actions.
\item
  The policy is sensitive to initialization when there are locally optimal actions close to initialization.
\end{enumerate}

Discretizing the action space or use Beta distribution helps avoid failure mode 1\&3 associated with Gaussian policy. Using KL regularization (same motivation as in \hyperref[trpo]{TRPO}) as an alternative surrogate model helps resolve failure mode 1\&2.

\begin{minipage}[t]{6.382978723404256e-3\linewidth}
\includegraphics{ppo-loss-functions.png}
\end{minipage}
\hfill
\begin{minipage}[t]{6.382978723404256e-3\linewidth}
\subsection{\texorpdfstring{PPG\hyperref[ppg]{\#}}{PPG\#}}\label{ppg}
\end{minipage}
\hfill
\begin{minipage}[t]{6.382978723404256e-3\linewidth}
{[}\href{https://arxiv.org/abs/2009.04416}{paper}\textbar{}\href{https://github.com/openai/phasic-policy-gradient}{code}{]}
\end{minipage}
\hfill
\begin{minipage}[t]{6.382978723404256e-3\linewidth}
Sharing parameters between policy and value networks have pros and cons. It allows policy and value functions to share the learned features with each other, but it may cause conflicts between competing objectives and demands the same data for training two networks at the same time. \textbf{Phasic policy gradient} (\textbf{PPG}; \href{https://arxiv.org/abs/2009.04416}{Cobbe, et al 2020}) modifies the traditional on-policy \hyperref[actor-critic]{actor-critic} policy gradient algorithm. precisely \hyperref[ppo]{PPO}, to have separate training phases for policy and value functions. In two alternating phases:
\end{minipage}
\hfill
\begin{minipage}[t]{6.382978723404256e-3\linewidth}
\begin{enumerate}
\tightlist
\item
  The \emph{policy phase}: updates the policy network by optimizing the PPO \hyperref[ppo_loss]{objective} \$L\^{}\textbackslash text\{CLIP\} (\textbackslash theta)\$;
\item
  The \emph{auxiliary phase}: optimizes an auxiliary objective alongside a behavioral cloning loss. In the paper, value function error is the sole auxiliary objective, but it can be quite general and includes any other additional auxiliary losses.
\end{enumerate}
\end{minipage}
\hfill
\begin{minipage}[t]{6.382978723404256e-3\linewidth}
\$\$ \textbackslash begin\{aligned\} L\^{}\textbackslash text\{joint\} \&= L\^{}\textbackslash text\{aux\} + \textbackslash beta\_\textbackslash text\{clone\} \textbackslash cdot \textbackslash mathbb\{E\}\_t{[}\textbackslash text\{KL\}{[}\textbackslash pi\_\{\textbackslash theta\_\textbackslash text\{old\}\}(\textbackslash cdot\textbackslash mid s\_t), \textbackslash pi\_\textbackslash theta(\textbackslash cdot\textbackslash mid s\_t){]}{]} \textbackslash\textbackslash{} L\^{}\textbackslash text\{aux\} \&= L\^{}\textbackslash text\{value\} = \textbackslash mathbb\{E\}\_t \textbackslash big{[}\textbackslash frac\{1\}\{2\}\textbackslash big( V\_w(s\_t) - \textbackslash hat\{V\}\_t\^{}\textbackslash text\{targ\} \textbackslash big)\^{}2\textbackslash big{]} \textbackslash end\{aligned\} \$\$
\end{minipage}
\hfill
\begin{minipage}[t]{6.382978723404256e-3\linewidth}
where \$\textbackslash beta\_\textbackslash text\{clone\}\$ is a hyperparameter for controlling how much we would like to keep the policy not diverge too much from its original behavior while optimizing the auxiliary objectives.
\end{minipage}
\hfill
\begin{subfigure}[t]{6.382978723404256e-3\linewidth}
\centering
\includegraphics{PPG_algo.png}
\caption{The algorithm of PPG. (Image source: \href{https://arxiv.org/abs/2009.04416}{Cobbe, et al 2020})}
\end{subfigure}
\hfill
\begin{minipage}[t]{6.382978723404256e-3\linewidth}
where
\end{minipage}
\hfill
\begin{minipage}[t]{6.382978723404256e-3\linewidth}
\begin{itemize}
\tightlist
\item
  \$N\_\textbackslash pi\$ is the number of policy update iterations in the policy phase. Note that the policy phase performs multiple iterations of updates per single auxiliary phase.
\item
  \$E\_\textbackslash pi\$ and \$E\_V\$ control the sample reuse (i.e. the number of training epochs performed across data in the reply buffer) for the policy and value functions, respectively. Note that this happens within the policy phase and thus \$E\_V\$ affects the learning of true value function not the auxiliary value function.
\item
  \$E\_\textbackslash text\{aux\}\$ defines the sample reuse in the auxiliary phrase. In PPG, value function optimization can tolerate a much higher level sample reuse; for example, in the experiments of the paper, \$E\_\textbackslash text\{aux\} = 6\$ while \$E\_\textbackslash pi = E\_V = 1\$.
\end{itemize}
\end{minipage}
\hfill
\begin{minipage}[t]{6.382978723404256e-3\linewidth}
PPG leads to a significant improvement on sample efficiency compared to PPO.
\end{minipage}
\hfill
\begin{subfigure}[t]{6.382978723404256e-3\linewidth}
\centering
\includegraphics{PPG_exp.png}
\caption{The mean normalized performance of PPG vs PPO on the \href{https://arxiv.org/abs/1912.01588}{Procgen} benchmark. (Image source: \href{https://arxiv.org/abs/2009.04416}{Cobbe, et al 2020})}
\end{subfigure}
\hfill
\begin{minipage}[t]{6.382978723404256e-3\linewidth}
\subsection{\texorpdfstring{ACER\hyperref[acer]{\#}}{ACER\#}}\label{acer}
\end{minipage}
\hfill
\begin{minipage}[t]{6.382978723404256e-3\linewidth}
{[}\href{https://arxiv.org/pdf/1611.01224.pdf}{paper}\textbar{}\href{https://github.com/openai/baselines/tree/master/baselines/acer}{code}{]}
\end{minipage}
\hfill
\begin{minipage}[t]{6.382978723404256e-3\linewidth}
\textbf{ACER}, short for \textbf{actor-critic with experience replay} (\href{https://arxiv.org/pdf/1611.01224.pdf}{Wang, et al., 2017}), is an off-policy actor-critic model with experience replay, greatly increasing the sample efficiency and decreasing the data correlation. A3C builds up the foundation for ACER, but it is on policy; ACER is A3C's off-policy counterpart. The major obstacle to making A3C off policy is how to control the stability of the off-policy estimator. ACER proposes three designs to overcome it:
\end{minipage}
\hfill
\begin{minipage}[t]{6.382978723404256e-3\linewidth}
\begin{itemize}
\tightlist
\item
  Use Retrace Q-value estimation;
\item
  Truncate the importance weights with bias correction;
\item
  Apply efficient TRPO.
\end{itemize}
\end{minipage}
\hfill
\begin{minipage}[t]{6.382978723404256e-3\linewidth}
\textbf{Retrace Q-value Estimation}
\end{minipage}
\hfill
\begin{minipage}[t]{6.382978723404256e-3\linewidth}
\href{http://papers.nips.cc/paper/6538-safe-and-efficient-off-policy-reinforcement-learning.pdf}{\emph{Retrace}} is an off-policy return-based Q-value estimation algorithm with a nice guarantee for convergence for any target and behavior policy pair \$(\textbackslash pi, \textbackslash beta)\$, plus good data efficiency.
\end{minipage}
\hfill
\begin{minipage}[t]{6.382978723404256e-3\linewidth}
Recall how TD learning works for prediction:
\end{minipage}
\hfill
\begin{minipage}[t]{6.382978723404256e-3\linewidth}
\begin{enumerate}
\tightlist
\item
  Compute TD error: \$\textbackslash delta\_t = R\_t + \textbackslash gamma \textbackslash mathbb\{E\}\_\{a \textbackslash sim \textbackslash pi\} Q(S\_\{t+1\}, a) - Q(S\_t, A\_t)\$; the term \$r\_t + \textbackslash gamma \textbackslash mathbb\{E\}\_\{a \textbackslash sim \textbackslash pi\} Q(s\_\{t+1\}, a) \$ is known as ``TD target''. The expectation \$\textbackslash mathbb\{E\}\_\{a \textbackslash sim \textbackslash pi\}\$ is used because for the future step the best estimation we can make is what the return would be if we follow the current policy \$\textbackslash pi\$.
\item
  Update the value by correcting the error to move toward the goal: \$Q(S\_t, A\_t) \textbackslash leftarrow Q(S\_t, A\_t) + \textbackslash alpha \textbackslash delta\_t\$. In other words, the incremental update on Q is proportional to the TD error: \$\textbackslash Delta Q(S\_t, A\_t) = \textbackslash alpha \textbackslash delta\_t\$.
\end{enumerate}
\end{minipage}
\hfill
\begin{minipage}[t]{6.382978723404256e-3\linewidth}
When the rollout is off policy, we need to apply importance sampling on the Q update:
\end{minipage}
\hfill
\begin{minipage}[t]{6.382978723404256e-3\linewidth}
\$\$ \textbackslash Delta Q\^{}\textbackslash text\{imp\}(S\_t, A\_t) = \textbackslash gamma\^{}t \textbackslash prod\_\{1 \textbackslash leq \textbackslash tau \textbackslash leq t\} \textbackslash frac\{\textbackslash pi(A\_\textbackslash tau \textbackslash vert S\_\textbackslash tau)\}\{\textbackslash beta(A\_\textbackslash tau \textbackslash vert S\_\textbackslash tau)\} \textbackslash delta\_t \$\$
\end{minipage}
\hfill
\begin{minipage}[t]{6.382978723404256e-3\linewidth}
The product of importance weights looks pretty scary when we start imagining how it can cause super high variance and even explode. Retrace Q-value estimation method modifies \$\textbackslash Delta Q\$ to have importance weights truncated by no more than a constant \$c\$:
\end{minipage}
\hfill
\begin{minipage}[t]{6.382978723404256e-3\linewidth}
\$\$ \textbackslash Delta Q\^{}\textbackslash text\{ret\}(S\_t, A\_t) = \textbackslash gamma\^{}t \textbackslash prod\_\{1 \textbackslash leq \textbackslash tau \textbackslash leq t\} \textbackslash min(c, \textbackslash frac\{\textbackslash pi(A\_\textbackslash tau \textbackslash vert S\_\textbackslash tau)\}\{\textbackslash beta(A\_\textbackslash tau \textbackslash vert S\_\textbackslash tau)\}) \textbackslash delta\_t \$\$
\end{minipage}
\hfill
\begin{minipage}[t]{6.382978723404256e-3\linewidth}
ACER uses \$Q\^{}\textbackslash text\{ret\}\$ as the target to train the critic by minimizing the L2 error term: \$(Q\^{}\textbackslash text\{ret\}(s, a) - Q(s, a))\^{}2\$.
\end{minipage}
\hfill
\begin{minipage}[t]{6.382978723404256e-3\linewidth}
\textbf{Importance weights truncation}
\end{minipage}
\hfill
\begin{minipage}[t]{6.382978723404256e-3\linewidth}
To reduce the high variance of the policy gradient \$\textbackslash hat\{g\}\$, ACER truncates the importance weights by a constant c, plus a correction term. The label \$\textbackslash hat\{g\}\_t\^{}\textbackslash text\{acer\}\$ is the ACER policy gradient at time t.
\end{minipage}
\hfill
\begin{minipage}[t]{6.382978723404256e-3\linewidth}
\$\$ \textbackslash begin\{aligned\} \textbackslash hat\{g\}\_t\^{}\textbackslash text\{acer\} = \& \textbackslash omega\_t \textbackslash big( Q\^{}\textbackslash text\{ret\}(S\_t, A\_t) - V\_\{\textbackslash theta\_v\}(S\_t) \textbackslash big) \textbackslash nabla\_\textbackslash theta \textbackslash ln \textbackslash pi\_\textbackslash theta(A\_t \textbackslash vert S\_t) \& \textbackslash scriptstyle\{\textbackslash text\{; Let \}\textbackslash omega\_t=\textbackslash frac\{\textbackslash pi(A\_t \textbackslash vert S\_t)\}\{\textbackslash beta(A\_t \textbackslash vert S\_t)\}\} \textbackslash\textbackslash{} = \& \textbackslash color\{blue\}\{\textbackslash min(c, \textbackslash omega\_t) \textbackslash big( Q\^{}\textbackslash text\{ret\}(S\_t, A\_t) - V\_w(S\_t) \textbackslash big) \textbackslash nabla\_\textbackslash theta \textbackslash ln \textbackslash pi\_\textbackslash theta(A\_t \textbackslash vert S\_t)\} \textbackslash\textbackslash{} \& + \textbackslash color\{red\}\{\textbackslash mathbb\{E\}\_\{a \textbackslash sim \textbackslash pi\} \textbackslash big{[} \textbackslash max(0, \textbackslash frac\{\textbackslash omega\_t(a) - c\}\{\textbackslash omega\_t(a)\}) \textbackslash big( Q\_w(S\_t, a) - V\_w(S\_t) \textbackslash big) \textbackslash nabla\_\textbackslash theta \textbackslash ln \textbackslash pi\_\textbackslash theta(a \textbackslash vert S\_t) \textbackslash big{]}\} \& \textbackslash scriptstyle\{\textbackslash text\{; Let \}\textbackslash omega\_t (a) =\textbackslash frac\{\textbackslash pi(a \textbackslash vert S\_t)\}\{\textbackslash beta(a \textbackslash vert S\_t)\}\} \textbackslash end\{aligned\} \$\$
\end{minipage}
\hfill
\begin{minipage}[t]{6.382978723404256e-3\linewidth}
where \$Q\_w(.)\$ and \$V\_w(.)\$ are value functions predicted by the critic with parameter w. The first term (blue) contains the clipped important weight. The clipping helps reduce the variance, in addition to subtracting state value function \$V\_w(.)\$ as a baseline. The second term (red) makes a correction to achieve unbiased estimation.
\end{minipage}
\hfill
\begin{minipage}[t]{6.382978723404256e-3\linewidth}
\textbf{Efficient TRPO}
\end{minipage}
\hfill
\begin{minipage}[t]{6.382978723404256e-3\linewidth}
Furthermore, ACER adopts the idea of TRPO but with a small adjustment to make it more computationally efficient: rather than measuring the KL divergence between policies before and after one update, ACER maintains a running average of past policies and forces the updated policy to not deviate far from this average.
\end{minipage}
\hfill
\begin{minipage}[t]{6.382978723404256e-3\linewidth}
The ACER \href{https://arxiv.org/pdf/1611.01224.pdf}{paper} is pretty dense with many equations. Hopefully, with the prior knowledge on TD learning, Q-learning, importance sampling and TRPO, you will find the \href{https://arxiv.org/pdf/1611.01224.pdf}{paper} slightly easier to follow :)
\end{minipage}
\hfill
\begin{minipage}[t]{6.382978723404256e-3\linewidth}
\subsection{\texorpdfstring{ACTKR\hyperref[actkr]{\#}}{ACTKR\#}}\label{actkr}
\end{minipage}
\hfill
\begin{minipage}[t]{6.382978723404256e-3\linewidth}
{[}\href{https://arxiv.org/pdf/1708.05144.pdf}{paper}\textbar{}\href{https://github.com/openai/baselines/tree/master/baselines/acktr}{code}{]}
\end{minipage}
\hfill
\begin{minipage}[t]{6.382978723404256e-3\linewidth}
\textbf{ACKTR (actor-critic using Kronecker-factored trust region)} (\href{https://arxiv.org/pdf/1708.05144.pdf}{Yuhuai Wu, et al., 2017}) proposed to use Kronecker-factored approximation curvature (\href{https://arxiv.org/pdf/1503.05671.pdf}{K-FAC}) to do the gradient update for both the critic and actor. K-FAC made an improvement on the computation of \emph{natural gradient}, which is quite different from our \emph{standard gradient}. \href{http://kvfrans.com/a-intuitive-explanation-of-natural-gradient-descent/}{Here} is a nice, intuitive explanation of natural gradient. One sentence summary is probably:
\end{minipage}
\hfill
\begin{minipage}[t]{6.382978723404256e-3\linewidth}
\begin{quote}
``we first consider all combinations of parameters that result in a new network a constant KL divergence away from the old network. This constant value can be viewed as the step size or learning rate. Out of all these possible combinations, we choose the one that minimizes our loss function.''
\end{quote}
\end{minipage}
\hfill
\begin{minipage}[t]{6.382978723404256e-3\linewidth}
I listed ACTKR here mainly for the completeness of this post, but I would not dive into details, as it involves a lot of theoretical knowledge on natural gradient and optimization methods. If interested, check these papers/posts, before reading the ACKTR paper:
\end{minipage}
\hfill
\begin{minipage}[t]{6.382978723404256e-3\linewidth}
\begin{itemize}
\tightlist
\item
  Amari. \href{http://citeseerx.ist.psu.edu/viewdoc/download?doi=10.1.1.452.7280&rep=rep1&type=pdf}{Natural Gradient Works Efficiently in Learning}. 1998
\item
  Kakade. \href{https://papers.nips.cc/paper/2073-a-natural-policy-gradient.pdf}{A Natural Policy Gradient}. 2002
\item
  \href{http://kvfrans.com/a-intuitive-explanation-of-natural-gradient-descent/}{A intuitive explanation of natural gradient descent}
\item
  \href{https://en.wikipedia.org/wiki/Kronecker_product}{Wiki: Kronecker product}
\item
  Martens \& Grosse. \href{http://proceedings.mlr.press/v37/martens15.pdf}{Optimizing neural networks with kronecker-factored approximate curvature.} 2015.
\end{itemize}
\end{minipage}
\hfill
\begin{minipage}[t]{6.382978723404256e-3\linewidth}
Here is a high level summary from the K-FAC \href{https://arxiv.org/pdf/1503.05671.pdf}{paper}:
\end{minipage}
\hfill
\begin{minipage}[t]{6.382978723404256e-3\linewidth}
\begin{quote}
``This approximation is built in two stages. In the first, the rows and columns of the Fisher are divided into groups, each of which corresponds to all the weights in a given layer, and this gives rise to a block-partitioning of the matrix. These blocks are then approximated as Kronecker products between much smaller matrices, which we show is equivalent to making certain approximating assumptions regarding the statistics of the network's gradients.
\end{quote}
\end{minipage}
\hfill
\begin{minipage}[t]{6.382978723404256e-3\linewidth}
\begin{quote}
In the second stage, this matrix is further approximated as having an inverse which is either block-diagonal or block-tridiagonal. We justify this approximation through a careful examination of the relationships between inverse covariances, tree-structured graphical models, and linear regression. Notably, this justification doesn't apply to the Fisher itself, and our experiments confirm that while the inverse Fisher does indeed possess this structure (approximately), the Fisher itself does not.''
\end{quote}
\end{minipage}
\hfill
\begin{minipage}[t]{6.382978723404256e-3\linewidth}
\subsection{\texorpdfstring{SAC\hyperref[sac]{\#}}{SAC\#}}\label{sac}
\end{minipage}
\hfill
\begin{minipage}[t]{6.382978723404256e-3\linewidth}
{[}\href{https://arxiv.org/abs/1801.01290}{paper}\textbar{}\href{https://github.com/haarnoja/sac}{code}{]}
\end{minipage}
\hfill
\begin{minipage}[t]{6.382978723404256e-3\linewidth}
\textbf{Soft Actor-Critic (SAC)} (\href{https://arxiv.org/abs/1801.01290}{Haarnoja et al. 2018}) incorporates the entropy measure of the policy into the reward to encourage exploration: we expect to learn a policy that acts as randomly as possible while it is still able to succeed at the task. It is an off-policy actor-critic model following the maximum entropy reinforcement learning framework. A precedent work is \href{https://arxiv.org/abs/1702.08165}{Soft Q-learning}.
\end{minipage}
\hfill
\begin{minipage}[t]{6.382978723404256e-3\linewidth}
Three key components in SAC:
\end{minipage}
\hfill
\begin{minipage}[t]{6.382978723404256e-3\linewidth}
\begin{itemize}
\tightlist
\item
  An \hyperref[actor-critic]{actor-critic} architecture with separate policy and value function networks;
\item
  An \hyperref[off-policy-policy-gradient]{off-policy} formulation that enables reuse of previously collected data for efficiency;
\item
  Entropy maximization to enable stability and exploration.
\end{itemize}
\end{minipage}
\hfill
\begin{minipage}[t]{6.382978723404256e-3\linewidth}
The policy is trained with the objective to maximize the expected return and the entropy at the same time:
\end{minipage}
\hfill
\begin{minipage}[t]{6.382978723404256e-3\linewidth}
\$\$ J(\textbackslash theta) = \textbackslash sum\_\{t=1\}\^{}T \textbackslash mathbb\{E\}\_\{(s\_t, a\_t) \textbackslash sim \textbackslash rho\_\{\textbackslash pi\_\textbackslash theta\}\} {[}r(s\_t, a\_t) + \textbackslash alpha \textbackslash mathcal\{H\}(\textbackslash pi\_\textbackslash theta(.\textbackslash vert s\_t)){]} \$\$
\end{minipage}
\hfill
\begin{minipage}[t]{6.382978723404256e-3\linewidth}
where \$\textbackslash mathcal\{H\}(.)\$ is the entropy measure and \$\textbackslash alpha\$ controls how important the entropy term is, known as \emph{temperature} parameter. The entropy maximization leads to policies that can (1) explore more and (2) capture multiple modes of near-optimal strategies (i.e., if there exist multiple options that seem to be equally good, the policy should assign each with an equal probability to be chosen).
\end{minipage}
\hfill
\begin{minipage}[t]{6.382978723404256e-3\linewidth}
Precisely, SAC aims to learn three functions:
\end{minipage}
\hfill
\begin{minipage}[t]{6.382978723404256e-3\linewidth}
\begin{itemize}
\tightlist
\item
  The policy with parameter \$\textbackslash theta\$, \$\textbackslash pi\_\textbackslash theta\$.
\item
  Soft Q-value function parameterized by \$w\$, \$Q\_w\$.
\item
  Soft state value function parameterized by \$\textbackslash psi\$, \$V\_\textbackslash psi\$; theoretically we can infer \$V\$ by knowing \$Q\$ and \$\textbackslash pi\$, but in practice, it helps stabilize the training.
\end{itemize}
\end{minipage}
\hfill
\begin{minipage}[t]{6.382978723404256e-3\linewidth}
Soft Q-value and soft state value are defined as:
\end{minipage}
\hfill
\begin{minipage}[t]{6.382978723404256e-3\linewidth}
\$\$ \textbackslash begin\{aligned\} Q(s\_t, a\_t) \&= r(s\_t, a\_t) + \textbackslash gamma \textbackslash mathbb\{E\}\_\{s\_\{t+1\} \textbackslash sim \textbackslash rho\_\{\textbackslash pi\}(s)\} {[}V(s\_\{t+1\}){]} \& \textbackslash text\{; according to Bellman equation.\}\textbackslash\textbackslash{} \textbackslash text\{where \}V(s\_t) \&= \textbackslash mathbb\{E\}\_\{a\_t \textbackslash sim \textbackslash pi\} {[}Q(s\_t, a\_t) - \textbackslash alpha \textbackslash log \textbackslash pi(a\_t \textbackslash vert s\_t){]} \& \textbackslash text\{; soft state value function.\} \textbackslash end\{aligned\} \$\$
\end{minipage}
\hfill
\begin{minipage}[t]{6.382978723404256e-3\linewidth}
\$\$ \textbackslash text\{Thus, \} Q(s\_t, a\_t) = r(s\_t, a\_t) + \textbackslash gamma \textbackslash mathbb\{E\}\_\{(s\_\{t+1\}, a\_\{t+1\}) \textbackslash sim \textbackslash rho\_\{\textbackslash pi\}\} {[}Q(s\_\{t+1\}, a\_\{t+1\}) - \textbackslash alpha \textbackslash log \textbackslash pi(a\_\{t+1\} \textbackslash vert s\_\{t+1\}){]} \$\$
\end{minipage}
\hfill
\begin{minipage}[t]{6.382978723404256e-3\linewidth}
\$\textbackslash rho\_\textbackslash pi(s)\$ and \$\textbackslash rho\_\textbackslash pi(s, a)\$ denote the state and the state-action marginals of the state distribution induced by the policy \$\textbackslash pi(a \textbackslash vert s)\$; see the similar definitions in \hyperref[dpg]{DPG} section.
\end{minipage}
\hfill
\begin{minipage}[t]{6.382978723404256e-3\linewidth}
The soft state value function is trained to minimize the mean squared error:
\end{minipage}
\hfill
\begin{minipage}[t]{6.382978723404256e-3\linewidth}
\$\$ \textbackslash begin\{aligned\} J\_V(\textbackslash psi) \&= \textbackslash mathbb\{E\}\_\{s\_t \textbackslash sim \textbackslash mathcal\{D\}\} {[}\textbackslash frac\{1\}\{2\} \textbackslash big(V\_\textbackslash psi(s\_t) - \textbackslash mathbb\{E\}{[}Q\_w(s\_t, a\_t) - \textbackslash log \textbackslash pi\_\textbackslash theta(a\_t \textbackslash vert s\_t){]} \textbackslash big)\^{}2{]} \textbackslash\textbackslash{} \textbackslash text\{with gradient: \}\textbackslash nabla\_\textbackslash psi J\_V(\textbackslash psi) \&= \textbackslash nabla\_\textbackslash psi V\_\textbackslash psi(s\_t)\textbackslash big( V\_\textbackslash psi(s\_t) - Q\_w(s\_t, a\_t) + \textbackslash log \textbackslash pi\_\textbackslash theta (a\_t \textbackslash vert s\_t) \textbackslash big) \textbackslash end\{aligned\} \$\$
\end{minipage}
\hfill
\begin{minipage}[t]{6.382978723404256e-3\linewidth}
where \$\textbackslash mathcal\{D\}\$ is the replay buffer.
\end{minipage}
\hfill
\begin{minipage}[t]{6.382978723404256e-3\linewidth}
The soft Q function is trained to minimize the soft Bellman residual:
\end{minipage}
\hfill
\begin{minipage}[t]{6.382978723404256e-3\linewidth}
\$\$ \textbackslash begin\{aligned\} J\_Q(w) \&= \textbackslash mathbb\{E\}\_\{(s\_t, a\_t) \textbackslash sim \textbackslash mathcal\{D\}\} {[}\textbackslash frac\{1\}\{2\}\textbackslash big( Q\_w(s\_t, a\_t) - (r(s\_t, a\_t) + \textbackslash gamma \textbackslash mathbb\{E\}\_\{s\_\{t+1\} \textbackslash sim \textbackslash rho\_\textbackslash pi(s)\}{[}V\_\{\textbackslash bar\{\textbackslash psi\}\}(s\_\{t+1\}){]}) \textbackslash big)\^{}2{]} \textbackslash\textbackslash{} \textbackslash text\{with gradient: \} \textbackslash nabla\_w J\_Q(w) \&= \textbackslash nabla\_w Q\_w(s\_t, a\_t) \textbackslash big( Q\_w(s\_t, a\_t) - r(s\_t, a\_t) - \textbackslash gamma V\_\{\textbackslash bar\{\textbackslash psi\}\}(s\_\{t+1\})\textbackslash big) \textbackslash end\{aligned\} \$\$
\end{minipage}
\hfill
\begin{minipage}[t]{6.382978723404256e-3\linewidth}
where \$\textbackslash bar\{\textbackslash psi\}\$ is the target value function which is the exponential moving average (or only gets updated periodically in a ``hard'' way), just like how the parameter of the target Q network is treated in \href{https://lilianweng.github.io/posts/2018-02-19-rl-overview/\#deep-q-network}{DQN} to stabilize the training.
\end{minipage}
\hfill
\begin{minipage}[t]{6.382978723404256e-3\linewidth}
SAC updates the policy to minimize the \href{https://en.wikipedia.org/wiki/Kullback\%E2\%80\%93Leibler_divergence}{KL-divergence}:
\end{minipage}
\hfill
\begin{minipage}[t]{6.382978723404256e-3\linewidth}
\$\$ \textbackslash begin\{aligned\} \textbackslash pi\_\textbackslash text\{new\} \&= \textbackslash arg\textbackslash min\_\{\textbackslash pi\textquotesingle{} \textbackslash in \textbackslash Pi\} D\_\textbackslash text\{KL\} \textbackslash Big( \textbackslash pi\textquotesingle(.\textbackslash vert s\_t) \textbackslash\textbar{} \textbackslash frac\{\textbackslash exp(Q\^{}\{\textbackslash pi\_\textbackslash text\{old\}\}(s\_t, .))\}\{Z\^{}\{\textbackslash pi\_\textbackslash text\{old\}\}(s\_t)\} \textbackslash Big) \textbackslash\textbackslash{[}6pt{]} \&= \textbackslash arg\textbackslash min\_\{\textbackslash pi\textquotesingle{} \textbackslash in \textbackslash Pi\} D\_\textbackslash text\{KL\} \textbackslash big( \textbackslash pi\textquotesingle(.\textbackslash vert s\_t) \textbackslash\textbar{} \textbackslash exp(Q\^{}\{\textbackslash pi\_\textbackslash text\{old\}\}(s\_t, .) - \textbackslash log Z\^{}\{\textbackslash pi\_\textbackslash text\{old\}\}(s\_t)) \textbackslash big) \textbackslash\textbackslash{[}6pt{]} \textbackslash text\{objective for update: \} J\_\textbackslash pi(\textbackslash theta) \&= \textbackslash nabla\_\textbackslash theta D\_\textbackslash text\{KL\} \textbackslash big( \textbackslash pi\_\textbackslash theta(. \textbackslash vert s\_t) \textbackslash\textbar{} \textbackslash exp(Q\_w(s\_t, .) - \textbackslash log Z\_w(s\_t)) \textbackslash big) \textbackslash\textbackslash{[}6pt{]} \&= \textbackslash mathbb\{E\}\_\{a\_t\textbackslash sim\textbackslash pi\} \textbackslash Big{[} - \textbackslash log \textbackslash big( \textbackslash frac\{\textbackslash exp(Q\_w(s\_t, a\_t) - \textbackslash log Z\_w(s\_t))\}\{\textbackslash pi\_\textbackslash theta(a\_t \textbackslash vert s\_t)\} \textbackslash big) \textbackslash Big{]} \textbackslash\textbackslash{[}6pt{]} \&= \textbackslash mathbb\{E\}\_\{a\_t\textbackslash sim\textbackslash pi\} {[} \textbackslash log \textbackslash pi\_\textbackslash theta(a\_t \textbackslash vert s\_t) - Q\_w(s\_t, a\_t) + \textbackslash log Z\_w(s\_t) {]} \textbackslash end\{aligned\} \$\$
\end{minipage}
\hfill
\begin{minipage}[t]{6.382978723404256e-3\linewidth}
where \$\textbackslash Pi\$ is the set of potential policies that we can model our policy as to keep them tractable; for example, \$\textbackslash Pi\$ can be the family of Gaussian mixture distributions, expensive to model but highly expressive and still tractable. \$Z\^{}\{\textbackslash pi\_\textbackslash text\{old\}\}(s\_t)\$ is the partition function to normalize the distribution. It is usually intractable but does not contribute to the gradient. How to minimize \$J\_\textbackslash pi(\textbackslash theta)\$ depends our choice of \$\textbackslash Pi\$.
\end{minipage}
\hfill
\begin{minipage}[t]{6.382978723404256e-3\linewidth}
This update guarantees that \$Q\^{}\{\textbackslash pi\_\textbackslash text\{new\}\}(s\_t, a\_t) \textbackslash geq Q\^{}\{\textbackslash pi\_\textbackslash text\{old\}\}(s\_t, a\_t)\$, please check the proof on this lemma in the Appendix B.2 in the original \href{https://arxiv.org/abs/1801.01290}{paper}.
\end{minipage}
\hfill
\begin{minipage}[t]{6.382978723404256e-3\linewidth}
Once we have defined the objective functions and gradients for soft action-state value, soft state value and the policy network, the soft actor-critic algorithm is straightforward:
\end{minipage}
\hfill
\begin{subfigure}[t]{6.382978723404256e-3\linewidth}
\centering
\includegraphics{SAC_algo.png}
\caption{The soft actor-critic algorithm. (Image source: \href{https://arxiv.org/abs/1801.01290}{original paper})}
\end{subfigure}
\hfill
\begin{minipage}[t]{6.382978723404256e-3\linewidth}
\subsection{\texorpdfstring{SAC with Automatically Adjusted Temperature\hyperref[sac-with-automatically-adjusted-temperature]{\#}}{SAC with Automatically Adjusted Temperature\#}}\label{sac-with-automatically-adjusted-temperature}
\end{minipage}
\hfill
\begin{minipage}[t]{6.382978723404256e-3\linewidth}
{[}\href{https://arxiv.org/abs/1812.05905}{paper}\textbar{}\href{https://github.com/rail-berkeley/softlearning}{code}{]}
\end{minipage}
\hfill
\begin{minipage}[t]{6.382978723404256e-3\linewidth}
SAC is brittle with respect to the temperature parameter. Unfortunately it is difficult to adjust temperature, because the entropy can vary unpredictably both across tasks and during training as the policy becomes better. An improvement on SAC formulates a constrained optimization problem: while maximizing the expected return, the policy should satisfy a minimum entropy constraint:
\end{minipage}
\hfill
\begin{minipage}[t]{6.382978723404256e-3\linewidth}
\$\$ \textbackslash max\_\{\textbackslash pi\_0, \textbackslash dots, \textbackslash pi\_T\} \textbackslash mathbb\{E\} \textbackslash Big{[} \textbackslash sum\_\{t=0\}\^{}T r(s\_t, a\_t)\textbackslash Big{]} \textbackslash text\{s.t. \} \textbackslash forall t\textbackslash text\{, \} \textbackslash mathcal\{H\}(\textbackslash pi\_t) \textbackslash geq \textbackslash mathcal\{H\}\_0 \$\$
\end{minipage}
\hfill
\begin{minipage}[t]{6.382978723404256e-3\linewidth}
where \$\textbackslash mathcal\{H\}\_0\$ is a predefined minimum policy entropy threshold.
\end{minipage}
\hfill
\begin{minipage}[t]{6.382978723404256e-3\linewidth}
The expected return \$\textbackslash mathbb\{E\} \textbackslash Big{[} \textbackslash sum\_\{t=0\}\^{}T r(s\_t, a\_t)\textbackslash Big{]}\$ can be decomposed into a sum of rewards at all the time steps. Because the policy \$\textbackslash pi\_t\$ at time t has no effect on the policy at the earlier time step, \$\textbackslash pi\_\{t-1\}\$, we can maximize the return at different steps backward in time --- this is essentially \textbf{DP}.
\end{minipage}
\hfill
\begin{minipage}[t]{6.382978723404256e-3\linewidth}
\$\$ \textbackslash underbrace\{\textbackslash max\_\{\textbackslash pi\_0\} \textbackslash Big( \textbackslash mathbb\{E\}{[}r(s\_0, a\_0){]}+ \textbackslash underbrace\{\textbackslash max\_\{\textbackslash pi\_1\} \textbackslash Big(\textbackslash mathbb\{E\}{[}...{]} + \textbackslash underbrace\{\textbackslash max\_\{\textbackslash pi\_T\} \textbackslash mathbb\{E\}{[}r(s\_T, a\_T){]}\}\_\textbackslash text\{1st maximization\} \textbackslash Big)\}\_\textbackslash text\{second but last maximization\} \textbackslash Big)\}\_\textbackslash text\{last maximization\} \$\$
\end{minipage}
\hfill
\begin{minipage}[t]{6.382978723404256e-3\linewidth}
where we consider \$\textbackslash gamma=1\$.
\end{minipage}
\hfill
\begin{minipage}[t]{6.382978723404256e-3\linewidth}
So we start the optimization from the last timestep \$T\$:
\end{minipage}
\hfill
\begin{minipage}[t]{6.382978723404256e-3\linewidth}
\$\$ \textbackslash text\{maximize \} \textbackslash mathbb\{E\}\_\{(s\_T, a\_T) \textbackslash sim \textbackslash rho\_\{\textbackslash pi\}\} {[} r(s\_T, a\_T) {]} \textbackslash text\{ s.t. \} \textbackslash mathcal\{H\}(\textbackslash pi\_T) - \textbackslash mathcal\{H\}\_0 \textbackslash geq 0 \$\$
\end{minipage}
\hfill
\begin{minipage}[t]{6.382978723404256e-3\linewidth}
First, let us define the following functions:
\end{minipage}
\hfill
\begin{minipage}[t]{6.382978723404256e-3\linewidth}
\$\$ \textbackslash begin\{aligned\} h(\textbackslash pi\_T) \&= \textbackslash mathcal\{H\}(\textbackslash pi\_T) - \textbackslash mathcal\{H\}\_0 = \textbackslash mathbb\{E\}\_\{(s\_T, a\_T) \textbackslash sim \textbackslash rho\_\{\textbackslash pi\}\} {[}-\textbackslash log \textbackslash pi\_T(a\_T\textbackslash vert s\_T){]} - \textbackslash mathcal\{H\}\_0\textbackslash\textbackslash{} f(\textbackslash pi\_T) \&= \textbackslash begin\{cases\} \textbackslash mathbb\{E\}\_\{(s\_T, a\_T) \textbackslash sim \textbackslash rho\_\{\textbackslash pi\}\} {[} r(s\_T, a\_T) {]}, \& \textbackslash text\{if \}h(\textbackslash pi\_T) \textbackslash geq 0 \textbackslash\textbackslash{} -\textbackslash infty, \& \textbackslash text\{otherwise\} \textbackslash end\{cases\} \textbackslash end\{aligned\} \$\$
\end{minipage}
\hfill
\begin{minipage}[t]{6.382978723404256e-3\linewidth}
And the optimization becomes:
\end{minipage}
\hfill
\begin{minipage}[t]{6.382978723404256e-3\linewidth}
\$\$ \textbackslash text\{maximize \} f(\textbackslash pi\_T) \textbackslash text\{ s.t. \} h(\textbackslash pi\_T) \textbackslash geq 0 \$\$
\end{minipage}
\hfill
\begin{minipage}[t]{6.382978723404256e-3\linewidth}
To solve the maximization optimization with inequality constraint, we can construct a \href{https://cs.stanford.edu/people/davidknowles/lagrangian_duality.pdf}{Lagrangian expression} with a Lagrange multiplier (also known as ``dual variable''), \$\textbackslash alpha\_T\$:
\end{minipage}
\hfill
\begin{minipage}[t]{6.382978723404256e-3\linewidth}
\$\$ L(\textbackslash pi\_T, \textbackslash alpha\_T) = f(\textbackslash pi\_T) + \textbackslash alpha\_T h(\textbackslash pi\_T) \$\$
\end{minipage}
\hfill
\begin{minipage}[t]{6.382978723404256e-3\linewidth}
Considering the case when we try to \emph{minimize \$L(\textbackslash pi\_T, \textbackslash alpha\_T)\$ with respect to \$\textbackslash alpha\_T\$} - given a particular value \$\textbackslash pi\_T\$,
\end{minipage}
\hfill
\begin{minipage}[t]{6.382978723404256e-3\linewidth}
\begin{itemize}
\tightlist
\item
  If the constraint is satisfied, \$h(\textbackslash pi\_T) \textbackslash geq 0\$, at best we can set \$\textbackslash alpha\_T=0\$ since we have no control over the value of \$f(\textbackslash pi\_T)\$. Thus, \$L(\textbackslash pi\_T, 0) = f(\textbackslash pi\_T)\$.
\item
  If the constraint is invalidated, \$h(\textbackslash pi\_T) \textless{} 0\$, we can achieve \$L(\textbackslash pi\_T, \textbackslash alpha\_T) \textbackslash to -\textbackslash infty\$ by taking \$\textbackslash alpha\_T \textbackslash to \textbackslash infty\$. Thus, \$L(\textbackslash pi\_T, \textbackslash infty) = -\textbackslash infty = f(\textbackslash pi\_T)\$.
\end{itemize}
\end{minipage}
\hfill
\begin{minipage}[t]{6.382978723404256e-3\linewidth}
In either case, we can recover the following equation,
\end{minipage}
\hfill
\begin{minipage}[t]{6.382978723404256e-3\linewidth}
\$\$ f(\textbackslash pi\_T) = \textbackslash min\_\{\textbackslash alpha\_T \textbackslash geq 0\} L(\textbackslash pi\_T, \textbackslash alpha\_T) \$\$
\end{minipage}
\hfill
\begin{minipage}[t]{6.382978723404256e-3\linewidth}
At the same time, we want to maximize \$f(\textbackslash pi\_T)\$,
\end{minipage}
\hfill
\begin{minipage}[t]{6.382978723404256e-3\linewidth}
\$\$ \textbackslash max\_\{\textbackslash pi\_T\} f(\textbackslash pi\_T) = \textbackslash min\_\{\textbackslash alpha\_T \textbackslash geq 0\} \textbackslash max\_\{\textbackslash pi\_T\} L(\textbackslash pi\_T, \textbackslash alpha\_T) \$\$
\end{minipage}
\hfill
\begin{minipage}[t]{6.382978723404256e-3\linewidth}
Therefore, to maximize \$f(\textbackslash pi\_T)\$, the dual problem is listed as below. Note that to make sure \$\textbackslash max\_\{\textbackslash pi\_T\} f(\textbackslash pi\_T)\$ is properly maximized and would not become \$-\textbackslash infty\$, the constraint has to be satisfied.
\end{minipage}
\hfill
\begin{minipage}[t]{6.382978723404256e-3\linewidth}
\$\$ \textbackslash begin\{aligned\} \textbackslash max\_\{\textbackslash pi\_T\} \textbackslash mathbb\{E\}{[} r(s\_T, a\_T) {]} \&= \textbackslash max\_\{\textbackslash pi\_T\} f(\textbackslash pi\_T) \textbackslash\textbackslash{} \&= \textbackslash min\_\{\textbackslash alpha\_T \textbackslash geq 0\} \textbackslash max\_\{\textbackslash pi\_T\} L(\textbackslash pi\_T, \textbackslash alpha\_T) \textbackslash\textbackslash{} \&= \textbackslash min\_\{\textbackslash alpha\_T \textbackslash geq 0\} \textbackslash max\_\{\textbackslash pi\_T\} f(\textbackslash pi\_T) + \textbackslash alpha\_T h(\textbackslash pi\_T) \textbackslash\textbackslash{} \&= \textbackslash min\_\{\textbackslash alpha\_T \textbackslash geq 0\} \textbackslash max\_\{\textbackslash pi\_T\} \textbackslash mathbb\{E\}\_\{(s\_T, a\_T) \textbackslash sim \textbackslash rho\_\{\textbackslash pi\}\} {[} r(s\_T, a\_T) {]} + \textbackslash alpha\_T ( \textbackslash mathbb\{E\}\_\{(s\_T, a\_T) \textbackslash sim \textbackslash rho\_\{\textbackslash pi\}\} {[}-\textbackslash log \textbackslash pi\_T(a\_T\textbackslash vert s\_T){]} - \textbackslash mathcal\{H\}\_0) \textbackslash\textbackslash{} \&= \textbackslash min\_\{\textbackslash alpha\_T \textbackslash geq 0\} \textbackslash max\_\{\textbackslash pi\_T\} \textbackslash mathbb\{E\}\_\{(s\_T, a\_T) \textbackslash sim \textbackslash rho\_\{\textbackslash pi\}\} {[} r(s\_T, a\_T) - \textbackslash alpha\_T \textbackslash log \textbackslash pi\_T(a\_T\textbackslash vert s\_T){]} - \textbackslash alpha\_T \textbackslash mathcal\{H\}\_0 \textbackslash\textbackslash{} \&= \textbackslash min\_\{\textbackslash alpha\_T \textbackslash geq 0\} \textbackslash max\_\{\textbackslash pi\_T\} \textbackslash mathbb\{E\}\_\{(s\_T, a\_T) \textbackslash sim \textbackslash rho\_\{\textbackslash pi\}\} {[} r(s\_T, a\_T) + \textbackslash alpha\_T \textbackslash mathcal\{H\}(\textbackslash pi\_T) - \textbackslash alpha\_T \textbackslash mathcal\{H\}\_0 {]} \textbackslash end\{aligned\} \$\$
\end{minipage}
\hfill
\begin{minipage}[t]{6.382978723404256e-3\linewidth}
We could compute the optimal \$\textbackslash pi\_T\$ and \$\textbackslash alpha\_T\$ iteratively. First given the current \$\textbackslash alpha\_T\$, get the best policy \$\textbackslash pi\_T\^{}\{*\}\$ that maximizes \$L(\textbackslash pi\_T\^{}\{*\}, \textbackslash alpha\_T)\$. Then plug in \$\textbackslash pi\_T\^{}\{*\}\$ and compute \$\textbackslash alpha\_T\^{}\{*\}\$ that minimizes \$L(\textbackslash pi\_T\^{}\{*\}, \textbackslash alpha\_T)\$. Assuming we have one neural network for policy and one network for temperature parameter, the iterative update process is more aligned with how we update network parameters during training.
\end{minipage}
\hfill
\begin{minipage}[t]{6.382978723404256e-3\linewidth}
\$\$ \textbackslash begin\{aligned\} \textbackslash pi\^{}\{*\}\_T \&= \textbackslash arg\textbackslash max\_\{\textbackslash pi\_T\} \textbackslash mathbb\{E\}\_\{(s\_T, a\_T) \textbackslash sim \textbackslash rho\_\{\textbackslash pi\}\} {[} r(s\_T, a\_T) + \textbackslash alpha\_T \textbackslash mathcal\{H\}(\textbackslash pi\_T) - \textbackslash alpha\_T \textbackslash mathcal\{H\}\_0 {]} \textbackslash\textbackslash{} \textbackslash color\{blue\}\{\textbackslash alpha\^{}\{*\}\_T\} \&\textbackslash color\{blue\}\{=\} \textbackslash color\{blue\}\{\textbackslash arg\textbackslash min\_\{\textbackslash alpha\_T \textbackslash geq 0\} \textbackslash mathbb\{E\}\_\{(s\_T, a\_T) \textbackslash sim \textbackslash rho\_\{\textbackslash pi\^{}\{*\}\}\} {[}\textbackslash alpha\_T \textbackslash mathcal\{H\}(\textbackslash pi\^{}\{*\}\_T) - \textbackslash alpha\_T \textbackslash mathcal\{H\}\_0 {]}\} \textbackslash end\{aligned\} \$\$
\end{minipage}
\hfill
\begin{minipage}[t]{6.382978723404256e-3\linewidth}
\$\$ \textbackslash text\{Thus, \}\textbackslash max\_\{\textbackslash pi\_T\} \textbackslash mathbb\{E\} {[} r(s\_T, a\_T) {]} = \textbackslash mathbb\{E\}\_\{(s\_T, a\_T) \textbackslash sim \textbackslash rho\_\{\textbackslash pi\^{}\{*\}\}\} {[} r(s\_T, a\_T) + \textbackslash alpha\^{}\{*\}\_T \textbackslash mathcal\{H\}(\textbackslash pi\^{}\{*\}\_T) - \textbackslash alpha\^{}\{*\}\_T \textbackslash mathcal\{H\}\_0 {]} \$\$
\end{minipage}
\hfill
\begin{minipage}[t]{6.382978723404256e-3\linewidth}
Now let's go back to the soft Q value function:
\end{minipage}
\hfill
\begin{minipage}[t]{6.382978723404256e-3\linewidth}
\$\$ \textbackslash begin\{aligned\} Q\_\{T-1\}(s\_\{T-1\}, a\_\{T-1\}) \&= r(s\_\{T-1\}, a\_\{T-1\}) + \textbackslash mathbb\{E\} {[}Q(s\_T, a\_T) - \textbackslash alpha\_T \textbackslash log \textbackslash pi(a\_T \textbackslash vert s\_T){]} \textbackslash\textbackslash{} \&= r(s\_\{T-1\}, a\_\{T-1\}) + \textbackslash mathbb\{E\} {[}r(s\_T, a\_T){]} + \textbackslash alpha\_T \textbackslash mathcal\{H\}(\textbackslash pi\_T) \textbackslash\textbackslash{} Q\_\{T-1\}\^{}\{*\}(s\_\{T-1\}, a\_\{T-1\}) \&= r(s\_\{T-1\}, a\_\{T-1\}) + \textbackslash max\_\{\textbackslash pi\_T\} \textbackslash mathbb\{E\} {[}r(s\_T, a\_T){]} + \textbackslash alpha\_T \textbackslash mathcal\{H\}(\textbackslash pi\^{}\{*\}\_T) \& \textbackslash text\{; plug in the optimal \}\textbackslash pi\_T\^{}\{*\} \textbackslash end\{aligned\} \$\$
\end{minipage}
\hfill
\begin{minipage}[t]{6.382978723404256e-3\linewidth}
Therefore the expected return is as follows, when we take one step further back to the time step \$T-1\$:
\end{minipage}
\hfill
\begin{minipage}[t]{6.382978723404256e-3\linewidth}
\$\$ \textbackslash begin\{aligned\} \&\textbackslash max\_\{\textbackslash pi\_\{T-1\}\}\textbackslash Big(\textbackslash mathbb\{E\}{[}r(s\_\{T-1\}, a\_\{T-1\}){]} + \textbackslash max\_\{\textbackslash pi\_T\} \textbackslash mathbb\{E\}{[}r(s\_T, a\_T{]} \textbackslash Big) \textbackslash\textbackslash{} \&= \textbackslash max\_\{\textbackslash pi\_\{T-1\}\} \textbackslash Big( Q\^{}\{*\}\_\{T-1\}(s\_\{T-1\}, a\_\{T-1\}) - \textbackslash alpha\^{}\{*\}\_T \textbackslash mathcal\{H\}(\textbackslash pi\^{}\{*\}\_T) \textbackslash Big) \& \textbackslash text\{; should s.t. \} \textbackslash mathcal\{H\}(\textbackslash pi\_\{T-1\}) - \textbackslash mathcal\{H\}\_0 \textbackslash geq 0 \textbackslash\textbackslash{} \&= \textbackslash min\_\{\textbackslash alpha\_\{T-1\} \textbackslash geq 0\} \textbackslash max\_\{\textbackslash pi\_\{T-1\}\} \textbackslash Big( Q\^{}\{*\}\_\{T-1\}(s\_\{T-1\}, a\_\{T-1\}) - \textbackslash alpha\^{}\{*\}\_T \textbackslash mathcal\{H\}(\textbackslash pi\^{}\{*\}\_T) + \textbackslash alpha\_\{T-1\} \textbackslash big( \textbackslash mathcal\{H\}(\textbackslash pi\_\{T-1\}) - \textbackslash mathcal\{H\}\_0 \textbackslash big) \textbackslash Big) \& \textbackslash text\{; dual problem w/ Lagrangian.\} \textbackslash\textbackslash{} \&= \textbackslash min\_\{\textbackslash alpha\_\{T-1\} \textbackslash geq 0\} \textbackslash max\_\{\textbackslash pi\_\{T-1\}\} \textbackslash Big( Q\^{}\{*\}\_\{T-1\}(s\_\{T-1\}, a\_\{T-1\}) + \textbackslash alpha\_\{T-1\} \textbackslash mathcal\{H\}(\textbackslash pi\_\{T-1\}) - \textbackslash alpha\_\{T-1\}\textbackslash mathcal\{H\}\_0 \textbackslash Big) - \textbackslash alpha\^{}\{*\}\_T \textbackslash mathcal\{H\}(\textbackslash pi\^{}\{*\}\_T) \textbackslash end\{aligned\} \$\$
\end{minipage}
\hfill
\begin{minipage}[t]{6.382978723404256e-3\linewidth}
Similar to the previous step,
\end{minipage}
\hfill
\begin{minipage}[t]{6.382978723404256e-3\linewidth}
\$\$ \textbackslash begin\{aligned\} \textbackslash pi\^{}\{*\}\_\{T-1\} \&= \textbackslash arg\textbackslash max\_\{\textbackslash pi\_\{T-1\}\} \textbackslash mathbb\{E\}\_\{(s\_\{T-1\}, a\_\{T-1\}) \textbackslash sim \textbackslash rho\_\textbackslash pi\} {[}Q\^{}\{*\}\_\{T-1\}(s\_\{T-1\}, a\_\{T-1\}) + \textbackslash alpha\_\{T-1\} \textbackslash mathcal\{H\}(\textbackslash pi\_\{T-1\}) - \textbackslash alpha\_\{T-1\} \textbackslash mathcal\{H\}\_0 {]} \textbackslash\textbackslash{} \textbackslash color\{green\}\{\textbackslash alpha\^{}\{*\}\_\{T-1\}\} \&\textbackslash color\{green\}\{=\} \textbackslash color\{green\}\{\textbackslash arg\textbackslash min\_\{\textbackslash alpha\_\{T-1\} \textbackslash geq 0\} \textbackslash mathbb\{E\}\_\{(s\_\{T-1\}, a\_\{T-1\}) \textbackslash sim \textbackslash rho\_\{\textbackslash pi\^{}\{*\}\}\} {[} \textbackslash alpha\_\{T-1\} \textbackslash mathcal\{H\}(\textbackslash pi\^{}\{*\}\_\{T-1\}) - \textbackslash alpha\_\{T-1\}\textbackslash mathcal\{H\}\_0 {]}\} \textbackslash end\{aligned\} \$\$
\end{minipage}
\hfill
\begin{minipage}[t]{6.382978723404256e-3\linewidth}
The equation for updating \$\textbackslash alpha\_\{T-1\}\$ in {green} has the same format as the equation for updating \$\textbackslash alpha\_\{T-1\}\$ in {blue} above. By repeating this process, we can learn the optimal temperature parameter in every step by minimizing the same objective function:
\end{minipage}
\hfill
\begin{minipage}[t]{6.382978723404256e-3\linewidth}
\$\$ J(\textbackslash alpha) = \textbackslash mathbb\{E\}\_\{a\_t \textbackslash sim \textbackslash pi\_t\} {[}-\textbackslash alpha \textbackslash log \textbackslash pi\_t(a\_t \textbackslash mid s\_t) - \textbackslash alpha \textbackslash mathcal\{H\}\_0{]} \$\$
\end{minipage}
\hfill
\begin{minipage}[t]{6.382978723404256e-3\linewidth}
The final algorithm is same as SAC except for learning \$\textbackslash alpha\$ explicitly with respect to the objective \$J(\textbackslash alpha)\$ (see Fig. 7):
\end{minipage}
\hfill
\begin{subfigure}[t]{6.382978723404256e-3\linewidth}
\centering
\includegraphics{SAC2_algo.png}
\caption{The soft actor-critic algorithm with automatically adjusted temperature. (Image source: \href{https://arxiv.org/abs/1812.05905}{original paper})}
\end{subfigure}
\hfill
\begin{minipage}[t]{6.382978723404256e-3\linewidth}
\subsection{\texorpdfstring{TD3\hyperref[td3]{\#}}{TD3\#}}\label{td3}
\end{minipage}
\hfill
\begin{minipage}[t]{6.382978723404256e-3\linewidth}
{[}\href{https://arxiv.org/abs/1802.09477}{paper}\textbar{}\href{https://github.com/sfujim/TD3}{code}{]}
\end{minipage}
\hfill
\begin{minipage}[t]{6.382978723404256e-3\linewidth}
The Q-learning algorithm is commonly known to suffer from the overestimation of the value function. This overestimation can propagate through the training iterations and negatively affect the policy. This property directly motivated \href{https://papers.nips.cc/paper/3964-double-q-learning}{Double Q-learning} and \href{https://arxiv.org/abs/1509.06461}{Double DQN}: the action selection and Q-value update are decoupled by using two value networks.
\end{minipage}
\hfill
\begin{minipage}[t]{6.382978723404256e-3\linewidth}
\textbf{Twin Delayed Deep Deterministic} (short for \textbf{TD3}; \href{https://arxiv.org/abs/1802.09477}{Fujimoto et al., 2018}) applied a couple of tricks on \hyperref[ddpg]{DDPG} to prevent the overestimation of the value function:
\end{minipage}
\hfill
\begin{minipage}[t]{6.382978723404256e-3\linewidth}
(1) \textbf{Clipped Double Q-learning}: In Double Q-Learning, the action selection and Q-value estimation are made by two networks separately. In the DDPG setting, given two deterministic actors \$(\textbackslash mu\_\{\textbackslash theta\_1\}, \textbackslash mu\_\{\textbackslash theta\_2\})\$ with two corresponding critics \$(Q\_\{w\_1\}, Q\_\{w\_2\})\$, the Double Q-learning Bellman targets look like:
\end{minipage}
\hfill
\begin{minipage}[t]{6.382978723404256e-3\linewidth}
\$\$ \textbackslash begin\{aligned\} y\_1 \&= r + \textbackslash gamma Q\_\{w\_2\}(s\textquotesingle, \textbackslash mu\_\{\textbackslash theta\_1\}(s\textquotesingle))\textbackslash\textbackslash{} y\_2 \&= r + \textbackslash gamma Q\_\{w\_1\}(s\textquotesingle, \textbackslash mu\_\{\textbackslash theta\_2\}(s\textquotesingle)) \textbackslash end\{aligned\} \$\$
\end{minipage}
\hfill
\begin{minipage}[t]{6.382978723404256e-3\linewidth}
However, due to the slow changing policy, these two networks could be too similar to make independent decisions. The \emph{Clipped Double Q-learning} instead uses the minimum estimation among two so as to favor underestimation bias which is hard to propagate through training:
\end{minipage}
\hfill
\begin{minipage}[t]{6.382978723404256e-3\linewidth}
\$\$ \textbackslash begin\{aligned\} y\_1 \&= r + \textbackslash gamma \textbackslash min\_\{i=1,2\}Q\_\{w\_i\}(s\textquotesingle, \textbackslash mu\_\{\textbackslash theta\_1\}(s\textquotesingle))\textbackslash\textbackslash{} y\_2 \&= r + \textbackslash gamma \textbackslash min\_\{i=1,2\} Q\_\{w\_i\}(s\textquotesingle, \textbackslash mu\_\{\textbackslash theta\_2\}(s\textquotesingle)) \textbackslash end\{aligned\} \$\$
\end{minipage}
\hfill
\begin{minipage}[t]{6.382978723404256e-3\linewidth}
(2) \textbf{Delayed update of Target and Policy Networks}: In the \href{https://lilianweng.github.io/posts/2018-02-19-rl-overview/\#actor-critic}{actor-critic} model, policy and value updates are deeply coupled: Value estimates diverge through overestimation when the policy is poor, and the policy will become poor if the value estimate itself is inaccurate.
\end{minipage}
\hfill
\begin{minipage}[t]{6.382978723404256e-3\linewidth}
To reduce the variance, TD3 updates the policy at a lower frequency than the Q-function. The policy network stays the same until the value error is small enough after several updates. The idea is similar to how the periodically-updated target network stay as a stable objective in \href{https://lilianweng.github.io/posts/2018-02-19-rl-overview/\#dqn}{DQN}.
\end{minipage}
\hfill
\begin{minipage}[t]{6.382978723404256e-3\linewidth}
(3) \textbf{Target Policy Smoothing}: Given a concern with deterministic policies that they can overfit to narrow peaks in the value function, TD3 introduced a smoothing regularization strategy on the value function: adding a small amount of clipped random noises to the selected action and averaging over mini-batches.
\end{minipage}
\hfill
\begin{minipage}[t]{6.382978723404256e-3\linewidth}
\$\$ \textbackslash begin\{aligned\} y \&= r + \textbackslash gamma Q\_w (s\textquotesingle, \textbackslash mu\_\{\textbackslash theta\}(s\textquotesingle) + \textbackslash epsilon) \& \textbackslash\textbackslash{} \textbackslash epsilon \&\textbackslash sim \textbackslash text\{clip\}(\textbackslash mathcal\{N\}(0, \textbackslash sigma), -c, +c) \& \textbackslash scriptstyle\{\textbackslash text\{ ; clipped random noises.\}\} \textbackslash end\{aligned\} \$\$
\end{minipage}
\hfill
\begin{minipage}[t]{6.382978723404256e-3\linewidth}
This approach mimics the idea of \href{https://lilianweng.github.io/posts/2018-02-19-rl-overview/\#sarsa-on-policy-td-control}{SARSA} update and enforces that similar actions should have similar values.
\end{minipage}
\hfill
\begin{minipage}[t]{6.382978723404256e-3\linewidth}
Here is the final algorithm:
\end{minipage}
\hfill
\begin{subfigure}[t]{6.382978723404256e-3\linewidth}
\centering
\includegraphics{TD3.png}
\caption{TD3 Algorithm. (Image source: \href{https://arxiv.org/abs/1802.09477}{Fujimoto et al., 2018})}
\end{subfigure}
\hfill
\begin{minipage}[t]{6.382978723404256e-3\linewidth}
\subsection{\texorpdfstring{SVPG\hyperref[svpg]{\#}}{SVPG\#}}\label{svpg}
\end{minipage}
\hfill
\begin{minipage}[t]{6.382978723404256e-3\linewidth}
{[}\href{https://arxiv.org/abs/1704.02399}{paper}\textbar{}\href{https://github.com/dilinwang820/Stein-Variational-Gradient-Descent}{code} for SVPG{]}
\end{minipage}
\hfill
\begin{minipage}[t]{6.382978723404256e-3\linewidth}
Stein Variational Policy Gradient (\textbf{SVPG}; \href{https://arxiv.org/abs/1704.02399}{Liu et al, 2017}) applies the \href{https://www.cs.dartmouth.edu/~qliu/stein.html}{Stein} variational gradient descent (\textbf{SVGD}; \href{https://arxiv.org/abs/1608.04471}{Liu and Wang, 2016}) algorithm to update the policy parameter \$\textbackslash theta\$.
\end{minipage}
\hfill
\begin{minipage}[t]{6.382978723404256e-3\linewidth}
In the setup of maximum entropy policy optimization, \$\textbackslash theta\$ is considered as a random variable \$\textbackslash theta \textbackslash sim q(\textbackslash theta)\$ and the model is expected to learn this distribution \$q(\textbackslash theta)\$. Assuming we know a prior on how \$q\$ might look like, \$q\_0\$, and we would like to guide the learning process to not make \$\textbackslash theta\$ too far away from \$q\_0\$ by optimizing the following objective function:
\end{minipage}
\hfill
\begin{minipage}[t]{6.382978723404256e-3\linewidth}
\$\$ \textbackslash hat\{J\}(\textbackslash theta) = \textbackslash mathbb\{E\}\_\{\textbackslash theta \textbackslash sim q\} {[}J(\textbackslash theta){]} - \textbackslash alpha D\_\textbackslash text\{KL\}(q\textbackslash\textbar q\_0) \$\$
\end{minipage}
\hfill
\begin{minipage}[t]{6.382978723404256e-3\linewidth}
where \$\textbackslash mathbb\{E\}\_\{\textbackslash theta \textbackslash sim q\} {[}R(\textbackslash theta){]}\$ is the expected reward when \$\textbackslash theta \textbackslash sim q(\textbackslash theta)\$ and \$D\_\textbackslash text\{KL\}\$ is the KL divergence.
\end{minipage}
\hfill
\begin{minipage}[t]{6.382978723404256e-3\linewidth}
If we don't have any prior information, we might set \$q\_0\$ as a uniform distribution and set \$q\_0(\textbackslash theta)\$ to a constant. Then the above objective function becomes \hyperref[SAC]{SAC}, where the entropy term encourages exploration:
\end{minipage}
\hfill
\begin{minipage}[t]{6.382978723404256e-3\linewidth}
\$\$ \textbackslash begin\{aligned\} \textbackslash hat\{J\}(\textbackslash theta) \&= \textbackslash mathbb\{E\}\_\{\textbackslash theta \textbackslash sim q\} {[}J(\textbackslash theta){]} - \textbackslash alpha D\_\textbackslash text\{KL\}(q\textbackslash\textbar q\_0) \textbackslash\textbackslash{} \&= \textbackslash mathbb\{E\}\_\{\textbackslash theta \textbackslash sim q\} {[}J(\textbackslash theta){]} - \textbackslash alpha \textbackslash mathbb\{E\}\_\{\textbackslash theta \textbackslash sim q\} {[}\textbackslash log q(\textbackslash theta) - \textbackslash log q\_0(\textbackslash theta){]} \textbackslash\textbackslash{} \&= \textbackslash mathbb\{E\}\_\{\textbackslash theta \textbackslash sim q\} {[}J(\textbackslash theta){]} + \textbackslash alpha H(q(\textbackslash theta)) \textbackslash end\{aligned\} \$\$
\end{minipage}
\hfill
\begin{minipage}[t]{6.382978723404256e-3\linewidth}
Let's take the derivative of \$\textbackslash hat\{J\}(\textbackslash theta) = \textbackslash mathbb\{E\}\_\{\textbackslash theta \textbackslash sim q\} {[}J(\textbackslash theta){]} - \textbackslash alpha D\_\textbackslash text\{KL\}(q\textbar q\_0)\$ w.r.t. \$q\$:
\end{minipage}
\hfill
\begin{minipage}[t]{6.382978723404256e-3\linewidth}
\$\$ \textbackslash begin\{aligned\} \textbackslash nabla\_q \textbackslash hat\{J\}(\textbackslash theta) \&= \textbackslash nabla\_q \textbackslash big( \textbackslash mathbb\{E\}\_\{\textbackslash theta \textbackslash sim q\} {[}J(\textbackslash theta){]} - \textbackslash alpha D\_\textbackslash text\{KL\}(q\textbackslash\textbar q\_0) \textbackslash big) \textbackslash\textbackslash{} \&= \textbackslash nabla\_q \textbackslash int\_\textbackslash theta \textbackslash big( q(\textbackslash theta) J(\textbackslash theta) - \textbackslash alpha q(\textbackslash theta)\textbackslash log q(\textbackslash theta) + \textbackslash alpha q(\textbackslash theta) \textbackslash log q\_0(\textbackslash theta) \textbackslash big) \textbackslash\textbackslash{} \&= \textbackslash int\_\textbackslash theta \textbackslash big( J(\textbackslash theta) - \textbackslash alpha \textbackslash log q(\textbackslash theta) -\textbackslash alpha + \textbackslash alpha \textbackslash log q\_0(\textbackslash theta) \textbackslash big) \textbackslash\textbackslash{} \&= 0 \textbackslash end\{aligned\} \$\$
\end{minipage}
\hfill
\begin{minipage}[t]{6.382978723404256e-3\linewidth}
The optimal distribution is:
\end{minipage}
\hfill
\begin{minipage}[t]{6.382978723404256e-3\linewidth}
\$\$ \textbackslash log q\^{}\{*\}(\textbackslash theta) = \textbackslash frac\{1\}\{\textbackslash alpha\} J(\textbackslash theta) + \textbackslash log q\_0(\textbackslash theta) - 1 \textbackslash text\{ thus \} \textbackslash underbrace\{ q\^{}\{*\}(\textbackslash theta) \}\_\textbackslash textrm\{"posterior"\} \textbackslash propto \textbackslash underbrace\{\textbackslash exp ( J(\textbackslash theta) / \textbackslash alpha )\}\_\textbackslash textrm\{"likelihood"\} \textbackslash underbrace\{q\_0(\textbackslash theta)\}\_\textbackslash textrm\{prior\} \$\$
\end{minipage}
\hfill
\begin{minipage}[t]{6.382978723404256e-3\linewidth}
The temperature \$\textbackslash alpha\$ decides a tradeoff between exploitation and exploration. When \$\textbackslash alpha \textbackslash rightarrow 0\$, \$\textbackslash theta\$ is updated only according to the expected return \$J(\textbackslash theta)\$. When \$\textbackslash alpha \textbackslash rightarrow \textbackslash infty\$, \$\textbackslash theta\$ always follows the prior belief.
\end{minipage}
\hfill
\begin{minipage}[t]{6.382978723404256e-3\linewidth}
When using the SVGD method to estimate the target posterior distribution \$q(\textbackslash theta)\$, it relies on a set of particle \$\textbackslash\{\textbackslash theta\_i\textbackslash\}\_\{i=1\}\^{}n\$ (independently trained policy agents) and each is updated:
\end{minipage}
\hfill
\begin{minipage}[t]{6.382978723404256e-3\linewidth}
\$\$ \textbackslash theta\_i \textbackslash gets \textbackslash theta\_i + \textbackslash epsilon \textbackslash phi\^{}\{*\}(\textbackslash theta\_i) \textbackslash text\{ where \} \textbackslash phi\^{}\{*\} = \textbackslash max\_\{\textbackslash phi \textbackslash in \textbackslash mathcal\{H\}\} \textbackslash\{ - \textbackslash nabla\_\textbackslash epsilon D\_\textbackslash text\{KL\} (q\textquotesingle\_\{{[}\textbackslash theta + \textbackslash epsilon \textbackslash phi(\textbackslash theta){]}\} \textbackslash\textbar{} q) \textbackslash text\{ s.t. \} \textbackslash\textbar\textbackslash phi\textbackslash\textbar\_\{\textbackslash mathcal\{H\}\} \textbackslash leq 1\textbackslash\} \$\$
\end{minipage}
\hfill
\begin{minipage}[t]{6.382978723404256e-3\linewidth}
where \$\textbackslash epsilon\$ is a learning rate and \$\textbackslash phi\^{}\{*\}\$ is the unit ball of a \href{http://mlss.tuebingen.mpg.de/2015/slides/gretton/part_1.pdf}{RKHS} (reproducing kernel Hilbert space) \$\textbackslash mathcal\{H\}\$ of \$\textbackslash theta\$-shaped value vectors that maximally decreases the KL divergence between the particles and the target distribution. \$q'(.)\$ is the distribution of \$\textbackslash theta + \textbackslash epsilon \textbackslash phi(\textbackslash theta)\$.
\end{minipage}
\hfill
\begin{minipage}[t]{6.382978723404256e-3\linewidth}
Comparing different gradient-based update methods:
\end{minipage}
\hfill
\begin{minipage}[t]{6.382978723404256e-3\linewidth}
\begin{longtable}[]{@{}ll@{}}
\toprule\noalign{}
Method & Update space \\
\midrule\noalign{}
\endhead
\bottomrule\noalign{}
\endlastfoot
Plain gradient & \$\textbackslash Delta \textbackslash theta\$ on the parameter space \\
\href{https://lilianweng.github.io/posts/2019-09-05-evolution-strategies/\#natural-gradients}{Natural gradient} & \$\textbackslash Delta \textbackslash theta\$ on the search distribution space \\
SVGD & \$\textbackslash Delta \textbackslash theta\$ on the kernel function space (edited) \\
\end{longtable}
\end{minipage}
\hfill
\begin{minipage}[t]{6.382978723404256e-3\linewidth}
One \href{https://arxiv.org/abs/1608.04471}{estimation} of \$\textbackslash phi\^{}\{*\}\$ has the following form. A positive definite kernel \$k(\textbackslash vartheta, \textbackslash theta)\$, i.e. a Gaussian \href{https://en.wikipedia.org/wiki/Radial_basis_function}{radial basis function}, measures the similarity between particles.
\end{minipage}
\hfill
\begin{minipage}[t]{6.382978723404256e-3\linewidth}
\$\$ \textbackslash begin\{aligned\} \textbackslash phi\^{}\{*\}(\textbackslash theta\_i) \&= \textbackslash mathbb\{E\}\_\{\textbackslash vartheta \textbackslash sim q\textquotesingle\} {[}\textbackslash nabla\_\textbackslash vartheta \textbackslash log q(\textbackslash vartheta) k(\textbackslash vartheta, \textbackslash theta\_i) + \textbackslash nabla\_\textbackslash vartheta k(\textbackslash vartheta, \textbackslash theta\_i){]}\textbackslash\textbackslash{} \&= \textbackslash frac\{1\}\{n\} \textbackslash sum\_\{j=1\}\^{}n {[}\textbackslash color\{red\}\{\textbackslash nabla\_\{\textbackslash theta\_j\} \textbackslash log q(\textbackslash theta\_j) k(\textbackslash theta\_j, \textbackslash theta\_i)\} + \textbackslash color\{green\}\{\textbackslash nabla\_\{\textbackslash theta\_j\} k(\textbackslash theta\_j, \textbackslash theta\_i)\}{]} \& \textbackslash scriptstyle\{\textbackslash text\{;approximate \}q\textquotesingle\textbackslash text\{ with current particle values\}\} \textbackslash end\{aligned\} \$\$
\end{minipage}
\hfill
\begin{minipage}[t]{6.382978723404256e-3\linewidth}
\begin{itemize}
\tightlist
\item
  The first term in {red} encourages \$\textbackslash theta\_i\$ learning towards the high probability regions of \$q\$ that is shared across similar particles. =\textgreater{} to be similar to other particles
\item
  The second term in {green} pushes particles away from each other and therefore diversifies the policy. =\textgreater{} to be dissimilar to other particles
\end{itemize}
\end{minipage}
\hfill
\begin{subfigure}[t]{6.382978723404256e-3\linewidth}
\centering
\begin{minipage}[t]{0.1285714285714286\linewidth}
\includegraphics{SVPG.png}
\end{minipage}
\hfill
\begin{minipage}[t]{0.1285714285714286\linewidth}
Usually the temperature \$\textbackslash alpha\$ follows an annealing scheme so that the training process does more exploration at the beginning but more exploitation at a later stage.
\end{minipage}
\hfill
\begin{minipage}[t]{0.1285714285714286\linewidth}
\subsection{\texorpdfstring{IMPALA\hyperref[impala]{\#}}{IMPALA\#}}\label{impala}
\end{minipage}
\hfill
\begin{minipage}[t]{0.1285714285714286\linewidth}
{[}\href{https://arxiv.org/abs/1802.01561}{paper}\textbar{}\href{https://github.com/deepmind/scalable_agent}{code}{]}
\end{minipage}
\hfill
\begin{minipage}[t]{0.1285714285714286\linewidth}
In order to scale up RL training to achieve a very high throughput, \textbf{IMPALA} (``Importance Weighted Actor-Learner Architecture'') framework decouples acting from learning on top of basic actor-critic setup and learns from all experience trajectories with \textbf{V-trace} off-policy correction.
\end{minipage}
\hfill
\begin{minipage}[t]{0.1285714285714286\linewidth}
Multiple actors generate experience in parallel, while the learner optimizes both policy and value function parameters using all the generated experience. Actors update their parameters with the latest policy from the learner periodically. Because acting and learning are decoupled, we can add many more actor machines to generate a lot more trajectories per time unit. As the training policy and the behavior policy are not totally synchronized, there is a \emph{gap} between them and thus we need off-policy corrections.
\end{minipage}
\hfill
\begin{subfigure}[t]{0.1285714285714286\linewidth}
\centering
\begin{minipage}[t]{6.923076923076923e-2\linewidth}
\includegraphics{IMPALA.png}
\end{minipage}
\hfill
\begin{minipage}[t]{6.923076923076923e-2\linewidth}
Let the value function \$V\_\textbackslash theta\$ parameterized by \$\textbackslash theta\$ and the policy \$\textbackslash pi\_\textbackslash phi\$ parameterized by \$\textbackslash phi\$. Also we know the trajectories in the replay buffer are collected by a slightly older policy \$\textbackslash mu\$.
\end{minipage}
\hfill
\begin{minipage}[t]{6.923076923076923e-2\linewidth}
At the training time \$t\$, given \$(s\_t, a\_t, s\_\{t+1\}, r\_t)\$, the value function parameter \$\textbackslash theta\$ is learned through an L2 loss between the current value and a V-trace value target. The \$n\$-step V-trace target is defined as:
\end{minipage}
\hfill
\begin{minipage}[t]{6.923076923076923e-2\linewidth}
\$\$ \textbackslash begin\{aligned\} v\_t \&= V\_\textbackslash theta(s\_t) + \textbackslash sum\_\{i=t\}\^{}\{t+n-1\} \textbackslash gamma\^{}\{i-t\} \textbackslash big(\textbackslash prod\_\{j=t\}\^{}\{i-1\} c\_j\textbackslash big) \textbackslash color\{red\}\{\textbackslash delta\_i V\} \textbackslash\textbackslash{} \&= V\_\textbackslash theta(s\_t) + \textbackslash sum\_\{i=t\}\^{}\{t+n-1\} \textbackslash gamma\^{}\{i-t\} \textbackslash big(\textbackslash prod\_\{j=t\}\^{}\{i-1\} c\_j\textbackslash big) \textbackslash color\{red\}\{\textbackslash rho\_i (r\_i + \textbackslash gamma V\_\textbackslash theta(s\_\{i+1\}) - V\_\textbackslash theta(s\_i))\} \textbackslash end\{aligned\} \$\$
\end{minipage}
\hfill
\begin{minipage}[t]{6.923076923076923e-2\linewidth}
where the {red} part \$\textbackslash delta\_i V\$ is a temporal difference for \$V\$. \$\textbackslash rho\_i = \textbackslash min\textbackslash big(\textbackslash bar\{\textbackslash rho\}, \textbackslash frac\{\textbackslash pi(a\_i \textbackslash vert s\_i)\}\{\textbackslash mu(a\_i \textbackslash vert s\_i)\}\textbackslash big)\$ and \$c\_j = \textbackslash min\textbackslash big(\textbackslash bar\{c\}, \textbackslash frac\{\textbackslash pi(a\_j \textbackslash vert s\_j)\}\{\textbackslash mu(a\_j \textbackslash vert s\_j)\}\textbackslash big)\$ are \emph{truncated \hyperref[off-policy-policy-gradient]{importance sampling (IS)} weights}. The product of \$c\_t, \textbackslash dots, c\_\{i-1\}\$ measures how much a temporal difference \$\textbackslash delta\_i V\$ observed at time \$i\$ impacts the update of the value function at a previous time \$t\$. In the on-policy case, we have \$\textbackslash rho\_i=1\$ and \$c\_j=1\$ (assuming \$\textbackslash bar\{c\} \textbackslash geq 1\$) and therefore the V-trace target becomes on-policy \$n\$-step Bellman target.
\end{minipage}
\hfill
\begin{minipage}[t]{6.923076923076923e-2\linewidth}
\$\textbackslash bar\{\textbackslash rho\}\$ and \$\textbackslash bar\{c\}\$ are two truncation constants with \$\textbackslash bar\{\textbackslash rho\} \textbackslash geq \textbackslash bar\{c\}\$. \$\textbackslash bar\{\textbackslash rho\}\$ impacts the fixed-point of the value function we converge to and \$\textbackslash bar\{c\}\$ impacts the speed of convergence. When \$\textbackslash bar\{\textbackslash rho\} =\textbackslash infty\$ (untruncated), we converge to the value function of the target policy \$V\^{}\textbackslash pi\$; when \$\textbackslash bar\{\textbackslash rho\}\$ is close to 0, we evaluate the value function of the behavior policy \$V\^{}\textbackslash mu\$; when in-between, we evaluate a policy between \$\textbackslash pi\$ and \$\textbackslash mu\$.
\end{minipage}
\hfill
\begin{minipage}[t]{6.923076923076923e-2\linewidth}
The value function parameter is therefore updated in the direction of:
\end{minipage}
\hfill
\begin{minipage}[t]{6.923076923076923e-2\linewidth}
\$\$ \textbackslash Delta\textbackslash theta = (v\_t - V\_\textbackslash theta(s\_t))\textbackslash nabla\_\textbackslash theta V\_\textbackslash theta(s\_t) \$\$
\end{minipage}
\hfill
\begin{minipage}[t]{6.923076923076923e-2\linewidth}
The policy parameter \$\textbackslash phi\$ is updated through policy gradient,
\end{minipage}
\hfill
\begin{minipage}[t]{6.923076923076923e-2\linewidth}
\$\$ \textbackslash begin\{aligned\} \textbackslash Delta \textbackslash phi \&= \textbackslash rho\_t \textbackslash nabla\_\textbackslash phi \textbackslash log \textbackslash pi\_\textbackslash phi(a\_t \textbackslash vert s\_t) \textbackslash big(r\_t + \textbackslash gamma v\_\{t+1\} - V\_\textbackslash theta(s\_t)\textbackslash big) + \textbackslash nabla\_\textbackslash phi H(\textbackslash pi\_\textbackslash phi)\textbackslash\textbackslash{} \&= \textbackslash rho\_t \textbackslash nabla\_\textbackslash phi \textbackslash log \textbackslash pi\_\textbackslash phi(a\_t \textbackslash vert s\_t) \textbackslash big(r\_t + \textbackslash gamma v\_\{t+1\} - V\_\textbackslash theta(s\_t)\textbackslash big) - \textbackslash nabla\_\textbackslash phi \textbackslash sum\_a \textbackslash pi\_\textbackslash phi(a\textbackslash vert s\_t)\textbackslash log \textbackslash pi\_\textbackslash phi(a\textbackslash vert s\_t) \textbackslash end\{aligned\} \$\$
\end{minipage}
\hfill
\begin{minipage}[t]{6.923076923076923e-2\linewidth}
where \$r\_t + \textbackslash gamma v\_\{t+1\}\$ is the estimated Q value, from which a state-dependent baseline \$V\_\textbackslash theta(s\_t)\$ is subtracted. \$H(\textbackslash pi\_\textbackslash phi)\$ is an entropy bonus to encourage exploration.
\end{minipage}
\hfill
\begin{minipage}[t]{6.923076923076923e-2\linewidth}
In the experiments, IMPALA is used to train one agent over multiple tasks. Two different model architectures are involved, a shallow model (left) and a deep residual model (right).
\end{minipage}
\hfill
\begin{subfigure}[t]{6.923076923076923e-2\linewidth}
\centering
\begin{minipage}[t]{2.25e-2\linewidth}
\includegraphics{IMPALA-arch.png}
\end{minipage}
\hfill
\begin{minipage}[t]{2.25e-2\linewidth}
\section{\texorpdfstring{Quick Summary\hyperref[quick-summary]{\#}}{Quick Summary\#}}\label{quick-summary}
\end{minipage}
\hfill
\begin{minipage}[t]{2.25e-2\linewidth}
After reading through all the algorithms above, I list a few building blocks or principles that seem to be common among them:
\end{minipage}
\hfill
\begin{minipage}[t]{2.25e-2\linewidth}
\begin{itemize}
\tightlist
\item
  Try to reduce the variance and keep the bias unchanged to stabilize learning.
\item
  Off-policy gives us better exploration and helps us use data samples more efficiently.
\item
  Experience replay (training data sampled from a replay memory buffer);
\item
  Target network that is either frozen periodically or updated slower than the actively learned policy network;
\item
  Batch normalization;
\item
  Entropy-regularized reward;
\item
  The critic and actor can share lower layer parameters of the network and two output heads for policy and value functions.
\item
  It is possible to learn with deterministic policy rather than stochastic one.
\item
  Put constraint on the divergence between policy updates.
\item
  New optimization methods (such as K-FAC).
\item
  Entropy maximization of the policy helps encourage exploration.
\item
  Try not to overestimate the value function.
\item
  Think twice whether the policy and value network should share parameters.
\item
  TBA more.
\end{itemize}
\end{minipage}
\hfill
\begin{minipage}[t]{2.25e-2\linewidth}
\begin{center}\rule{0.5\linewidth}{0.5pt}\end{center}
\end{minipage}
\hfill
\begin{minipage}[t]{2.25e-2\linewidth}
Cited as:
\end{minipage}
\hfill
\begin{minipage}[t]{2.25e-2\linewidth}
\begin{verbatim}
@article{weng2018PG,
  title   = "Policy Gradient Algorithms",
  author  = "Weng, Lilian",
  journal = "lilianweng.github.io",
  year    = "2018",
  url     = "https://lilianweng.github.io/posts/2018-04-08-policy-gradient/"
}
\end{verbatim}
\end{minipage}
\hfill
\begin{minipage}[t]{2.25e-2\linewidth}
\section{\texorpdfstring{References\hyperref[references]{\#}}{References\#}}\label{references}
\end{minipage}
\hfill
\begin{minipage}[t]{2.25e-2\linewidth}
{[}1{]} jeremykun.com \href{https://jeremykun.com/2015/04/06/markov-chain-monte-carlo-without-all-the-bullshit/}{Markov Chain Monte Carlo Without all the Bullshit}
\end{minipage}
\hfill
\begin{minipage}[t]{2.25e-2\linewidth}
{[}2{]} Richard S. Sutton and Andrew G. Barto. \href{http://incompleteideas.net/book/bookdraft2017nov5.pdf}{Reinforcement Learning: An Introduction; 2nd Edition}. 2017.
\end{minipage}
\hfill
\begin{minipage}[t]{2.25e-2\linewidth}
{[}3{]} John Schulman, et al. \href{https://arxiv.org/pdf/1506.02438.pdf}{``High-dimensional continuous control using generalized advantage estimation.''} ICLR 2016.
\end{minipage}
\hfill
\begin{minipage}[t]{2.25e-2\linewidth}
{[}4{]} Thomas Degris, Martha White, and Richard S. Sutton. \href{https://arxiv.org/pdf/1205.4839.pdf}{``Off-policy actor-critic.''} ICML 2012.
\end{minipage}
\hfill
\begin{minipage}[t]{2.25e-2\linewidth}
{[}5{]} timvieira.github.io \href{http://timvieira.github.io/blog/post/2014/12/21/importance-sampling/}{Importance sampling}
\end{minipage}
\hfill
\begin{minipage}[t]{2.25e-2\linewidth}
{[}6{]} Mnih, Volodymyr, et al. \href{https://arxiv.org/abs/1602.01783}{``Asynchronous methods for deep reinforcement learning.''} ICML. 2016.
\end{minipage}
\hfill
\begin{minipage}[t]{2.25e-2\linewidth}
{[}7{]} David Silver, et al. \href{https://hal.inria.fr/file/index/docid/938992/filename/dpg-icml2014.pdf}{``Deterministic policy gradient algorithms.''} ICML. 2014.
\end{minipage}
\hfill
\begin{minipage}[t]{2.25e-2\linewidth}
{[}8{]} Timothy P. Lillicrap, et al. \href{https://arxiv.org/pdf/1509.02971.pdf}{``Continuous control with deep reinforcement learning.''} arXiv preprint arXiv:1509.02971 (2015).
\end{minipage}
\hfill
\begin{minipage}[t]{2.25e-2\linewidth}
{[}9{]} Ryan Lowe, et al. \href{https://arxiv.org/pdf/1706.02275.pdf}{``Multi-agent actor-critic for mixed cooperative-competitive environments.''} NIPS. 2017.
\end{minipage}
\hfill
\begin{minipage}[t]{2.25e-2\linewidth}
{[}10{]} John Schulman, et al. \href{https://arxiv.org/pdf/1502.05477.pdf}{``Trust region policy optimization.''} ICML. 2015.
\end{minipage}
\hfill
\begin{minipage}[t]{2.25e-2\linewidth}
{[}11{]} Ziyu Wang, et al. \href{https://arxiv.org/pdf/1611.01224.pdf}{``Sample efficient actor-critic with experience replay.''} ICLR 2017.
\end{minipage}
\hfill
\begin{minipage}[t]{2.25e-2\linewidth}
{[}12{]} Rémi Munos, Tom Stepleton, Anna Harutyunyan, and Marc Bellemare. \href{http://papers.nips.cc/paper/6538-safe-and-efficient-off-policy-reinforcement-learning.pdf}{``Safe and efficient off-policy reinforcement learning''} NIPS. 2016.
\end{minipage}
\hfill
\begin{minipage}[t]{2.25e-2\linewidth}
{[}13{]} Yuhuai Wu, et al. \href{https://arxiv.org/pdf/1708.05144.pdf}{``Scalable trust-region method for deep reinforcement learning using Kronecker-factored approximation.''} NIPS. 2017.
\end{minipage}
\hfill
\begin{minipage}[t]{2.25e-2\linewidth}
{[}14{]} kvfrans.com \href{http://kvfrans.com/a-intuitive-explanation-of-natural-gradient-descent/}{A intuitive explanation of natural gradient descent}
\end{minipage}
\hfill
\begin{minipage}[t]{2.25e-2\linewidth}
{[}15{]} Sham Kakade. \href{https://papers.nips.cc/paper/2073-a-natural-policy-gradient.pdf}{``A Natural Policy Gradient.''}. NIPS. 2002.
\end{minipage}
\hfill
\begin{minipage}[t]{2.25e-2\linewidth}
{[}16{]} \href{https://danieltakeshi.github.io/2017/03/28/going-deeper-into-reinforcement-learning-fundamentals-of-policy-gradients/}{``Going Deeper Into Reinforcement Learning: Fundamentals of Policy Gradients.''} - Seita's Place, Mar 2017.
\end{minipage}
\hfill
\begin{minipage}[t]{2.25e-2\linewidth}
{[}17{]} \href{https://danieltakeshi.github.io/2017/04/02/notes-on-the-generalized-advantage-estimation-paper/}{``Notes on the Generalized Advantage Estimation Paper.''} - Seita's Place, Apr, 2017.
\end{minipage}
\hfill
\begin{minipage}[t]{2.25e-2\linewidth}
{[}18{]} Gabriel Barth-Maron, et al. \href{https://arxiv.org/pdf/1804.08617.pdf}{``Distributed Distributional Deterministic Policy Gradients.''} ICLR 2018 poster.
\end{minipage}
\hfill
\begin{minipage}[t]{2.25e-2\linewidth}
{[}19{]} Tuomas Haarnoja, Aurick Zhou, Pieter Abbeel, and Sergey Levine. \href{https://arxiv.org/pdf/1801.01290.pdf}{``Soft Actor-Critic: Off-Policy Maximum Entropy Deep Reinforcement Learning with a Stochastic Actor.''} arXiv preprint arXiv:1801.01290 (2018).
\end{minipage}
\hfill
\begin{minipage}[t]{2.25e-2\linewidth}
{[}20{]} Scott Fujimoto, Herke van Hoof, and Dave Meger. \href{https://arxiv.org/abs/1802.09477}{``Addressing Function Approximation Error in Actor-Critic Methods.''} arXiv preprint arXiv:1802.09477 (2018).
\end{minipage}
\hfill
\begin{minipage}[t]{2.25e-2\linewidth}
{[}21{]} Tuomas Haarnoja, et al. \href{https://arxiv.org/abs/1812.05905}{``Soft Actor-Critic Algorithms and Applications.''} arXiv preprint arXiv:1812.05905 (2018).
\end{minipage}
\hfill
\begin{minipage}[t]{2.25e-2\linewidth}
{[}22{]} David Knowles. \href{https://cs.stanford.edu/people/davidknowles/lagrangian_duality.pdf}{``Lagrangian Duality for Dummies''} Nov 13, 2010.
\end{minipage}
\hfill
\begin{minipage}[t]{2.25e-2\linewidth}
{[}23{]} Yang Liu, et al. \href{https://arxiv.org/abs/1704.02399}{``Stein variational policy gradient.''} arXiv preprint arXiv:1704.02399 (2017).
\end{minipage}
\hfill
\begin{minipage}[t]{2.25e-2\linewidth}
{[}24{]} Qiang Liu and Dilin Wang. \href{https://papers.nips.cc/paper/6338-stein-variational-gradient-descent-a-general-purpose-bayesian-inference-algorithm.pdf}{``Stein variational gradient descent: A general purpose bayesian inference algorithm.''} NIPS. 2016.
\end{minipage}
\hfill
\begin{minipage}[t]{2.25e-2\linewidth}
{[}25{]} Lasse Espeholt, et al. \href{https://arxiv.org/abs/1802.01561}{``IMPALA: Scalable Distributed Deep-RL with Importance Weighted Actor-Learner Architectures''} arXiv preprint 1802.01561 (2018).
\end{minipage}
\hfill
\begin{minipage}[t]{2.25e-2\linewidth}
{[}26{]} Karl Cobbe, et al. \href{https://arxiv.org/abs/2009.04416}{``Phasic Policy Gradient.''} arXiv preprint arXiv:2009.04416 (2020).
\end{minipage}
\hfill
\begin{minipage}[t]{2.25e-2\linewidth}
{[}27{]} Chloe Ching-Yun Hsu, et al. \href{https://arxiv.org/abs/2009.10897}{``Revisiting Design Choices in Proximal Policy Optimization.''} arXiv preprint arXiv:2009.10897 (2020).
\end{minipage}
\hfill
\begin{minipage}[t]{2.25e-2\linewidth}
\begin{itemize}
\tightlist
\item
  \href{https://lilianweng.github.io/tags/reinforcement-learning/}{Reinforcement-Learning}
\item
  \href{https://lilianweng.github.io/tags/long-read/}{Long-Read}
\item
  \href{https://lilianweng.github.io/tags/math-heavy/}{Math-Heavy}
\end{itemize}
\end{minipage}
\hfill
\begin{minipage}[t]{2.25e-2\linewidth}
\href{https://lilianweng.github.io/posts/2018-05-05-drl-implementation/}{{« }\\
{Implementing Deep Reinforcement Learning Models with Tensorflow + OpenAI Gym}} \href{https://lilianweng.github.io/posts/2018-02-19-rl-overview/}{{ »}\\
{A (Long) Peek into Reinforcement Learning}}
\end{minipage}
\hfill
\begin{minipage}[t]{2.25e-2\linewidth}
\href{https://twitter.com/intent/tweet/?text=Policy\%20Gradient\%20Algorithms&url=https\%3a\%2f\%2flilianweng.github.io\%2fposts\%2f2018-04-08-policy-gradient\%2f&hashtags=reinforcement-learning\%2clong-read\%2cmath-heavy}{} \href{https://www.linkedin.com/shareArticle?mini=true&url=https\%3a\%2f\%2flilianweng.github.io\%2fposts\%2f2018-04-08-policy-gradient\%2f&title=Policy\%20Gradient\%20Algorithms&summary=Policy\%20Gradient\%20Algorithms&source=https\%3a\%2f\%2flilianweng.github.io\%2fposts\%2f2018-04-08-policy-gradient\%2f}{} \href{https://reddit.com/submit?url=https\%3a\%2f\%2flilianweng.github.io\%2fposts\%2f2018-04-08-policy-gradient\%2f&title=Policy\%20Gradient\%20Algorithms}{} \href{https://facebook.com/sharer/sharer.php?u=https\%3a\%2f\%2flilianweng.github.io\%2fposts\%2f2018-04-08-policy-gradient\%2f}{} \href{https://api.whatsapp.com/send?text=Policy\%20Gradient\%20Algorithms\%20-\%20https\%3a\%2f\%2flilianweng.github.io\%2fposts\%2f2018-04-08-policy-gradient\%2f}{} \href{https://telegram.me/share/url?text=Policy\%20Gradient\%20Algorithms&url=https\%3a\%2f\%2flilianweng.github.io\%2fposts\%2f2018-04-08-policy-gradient\%2f}{}
\end{minipage}
\hfill
\begin{minipage}[t]{2.25e-2\linewidth}
{© 2025 \href{https://lilianweng.github.io/}{Lil\textquotesingle Log}} { Powered by \href{https://gohugo.io/}{Hugo} \& \href{https://git.io/hugopapermod}{PaperMod} }
\end{minipage}
\hfill
\begin{minipage}[t]{2.25e-2\linewidth}
\phantomsection\label{top-link}\hyperref[top]{}
\end{minipage}
\caption{}
\end{subfigure}
\caption{}
\end{subfigure}
\caption{}
\end{subfigure}
